% =============================================================================
% Recon-Platform-R2 — Lucrare de finalizare a studiilor (UPT, Facultatea AC)
% -----------------------------------------------------------------------------
% Șablon respectă "Template AC ro v2.docx":
%   - Arial 12 pt (aproximat prin Helvetica/Nimbus Sans via pachetul helvet)
%   - Spațiere între rânduri: 1,15
%   - Margini oglindite: interior 2 cm, exterior 2 cm, sus 2,5 cm, jos 2 cm
%   - Aliniere justified, primul rând al paragrafului 1,5 cm indent
%   - Heading 1 (capitol): 14 pt, Bold, Uppercase, Center
%   - Heading 2 (secțiune): 12 pt, Bold, Uppercase, Left
%   - Antet: începând cu Introducerea (UPT Timișoara + program + nume + titlu)
%   - Subsol: număr de pagină centrat, începând cu Introducerea
%   - Stil de citare: IEEE (numeric, ordinea apariției)
%   - Diacritice românești corecte: ă, â, î, ș (U+0219), ț (U+021B)
%
% Compilare:    pdflatex main → bibtex main → pdflatex main → pdflatex main
% (Overleaf compilează automat secvența cu setarea "Latexmk" + "pdfLaTeX".)
% =============================================================================

\documentclass[a4paper,12pt,twoside,openany]{report}

% --- Codificare și limbă -----------------------------------------------------
\usepackage[utf8]{inputenc}
\usepackage[T1]{fontenc}
% Latin Modern: fonturi vectoriale scalabile la orice dimensiune. Elimină
% avertismentele „Font shape ... size <16>/<20> not available, substituted"
% care apar la titlurile de 16/20 pt (matematica din heading-uri nu există în
% setul bitmap clasic Computer Modern). Textul rămâne Helvetica (vezi helvet).
\usepackage{lmodern}
% Ordinea contează: ultima limbă din listă devine cea principală.
% English este încărcat ca auxiliară pentru blocul ABSTRACT.
\usepackage[english,romanian]{babel}

% Font sans-serif similar cu Arial (Nimbus Sans / Helvetica) la dimensiunea de bază
\usepackage{helvet}
\renewcommand{\familydefault}{\sfdefault}

% microtype: protruziune + expansiune de font la justificare. Reduce
% considerabil liniile „Overfull \hbox" (text care iese în marginea dreaptă),
% fără gap-uri vizibile între cuvinte. Compatibil cu pdfLaTeX.
\usepackage{microtype}

% --- Margini (oglindite) -----------------------------------------------------
\usepackage[
    a4paper,
    twoside,
    inner=2cm,
    outer=2cm,
    top=2.5cm,
    bottom=2cm,
    headheight=1.5cm,
    headsep=0.5cm,
    footskip=1cm,
    includehead,
    includefoot
]{geometry}

% --- Spațiere între rânduri --------------------------------------------------
\usepackage{setspace}
\setstretch{1.15}

% Permite paginilor să nu fie umplute complet pe verticală — elimină
% avertismentele „Underfull \vbox (badness 10000) ... while \output is active"
% care apar firesc la paginile de capitol / \clearpage.
\raggedbottom
% Marjă suplimentară de întindere a rândurilor înainte de a declara o linie
% „Overfull" — reduce mult depășirile cauzate de identificatori lungi din
% \texttt{...} (nume de fișiere, apeluri de funcții) care nu se pot despărți.
\setlength{\emergencystretch}{4em}

% --- Indentare prim rând și paragraf ----------------------------------------
\usepackage{indentfirst}
\setlength{\parindent}{1.5cm}
\setlength{\parskip}{0pt}

% --- Antet / subsol ----------------------------------------------------------
\usepackage{fancyhdr}
\pagestyle{fancy}
\fancyhf{}
% Antet conform „Template AC ro v2.docx" regula (f): text de înălțime 8 pt,
% aliniat la stânga, pe rânduri succesive — (i) Universitatea Politehnica
% Timișoara; (ii) programul de studii + anul; (iii) numele candidatului;
% (iv) titlul lucrării. În dreapta, sigla UPT.
\fancyhead[L]{%
    \fontsize{8}{9.6}\selectfont
    Universitatea Politehnica Timișoara\\
    \specializare{} --- Anul universitar 2025--2026\\
    \candidatname{} --- Valentin-Gabriel NICOLA\\
    \workTitle{}%
}
% Sigla UPT (extrasă din „Template AC ro v2.docx", word/media/image2.jpeg).
% Fișierul upt_logo.jpeg trebuie încărcat alături de main.tex pe Overleaf.
\fancyhead[R]{\includegraphics[height=1.0cm]{upt_logo.jpeg}}
% Subsol: număr de pagină centrat, în formatul „- N -" (ca în footer1.xml).
\fancyfoot[C]{- \thepage{} -}
\renewcommand{\headrulewidth}{0.4pt}
\renewcommand{\footrulewidth}{0pt}

% Stil dedicat paginilor de titlu / rezumat / abstract / cuprins
% (fără antet, fără număr de pagină).
\fancypagestyle{coperta}{%
    \fancyhf{}%
    \renewcommand{\headrulewidth}{0pt}%
}
% Paginile de deschidere de capitol folosesc implicit stilul „plain" al clasei
% report. În front matter îl vrem GOL (ca \tableofcontents pe mai multe pagini
% să nu afișeze antet/număr); din Introducere îl aliem la „fancy" (mai jos), ca
% antetul + numărul să apară pe TOATE paginile, inclusiv cele de început de capitol.
\fancypagestyle{plain}{%
    \fancyhf{}%
    \renewcommand{\headrulewidth}{0pt}%
}

% --- Stiluri pentru titluri (per cerințele template-ului) -------------------
\usepackage{titlesec}

% Capitol = Heading 1: 14 pt, Bold, Uppercase, Center, numerotat „N. TITLU"
% pe un singur rând (forma [block] centrează linia; titlesec înlocuiește
% complet antetul de capitol al clasei report, deci cuvântul „Capitolul" NU apare).
\titleformat{\chapter}[block]
    {\centering\fontsize{14pt}{17pt}\bfseries\MakeUppercase}
    {\thechapter.}{0.5em}{}
\titlespacing*{\chapter}{0pt}{0pt}{24pt} % 2 rânduri libere de 12 pt sub titlu

% Secțiune = Heading 2: 12 pt, Bold, Uppercase, Left
\titleformat{\section}[hang]
    {\fontsize{12pt}{14pt}\bfseries\MakeUppercase}
    {\thesection}{0.5em}{}
\titlespacing*{\section}{0pt}{14pt}{12pt}

% Subsecțiune (opțional)
\titleformat{\subsection}[hang]
    {\fontsize{12pt}{14pt}\bfseries}
    {\thesubsection}{0.5em}{}

% --- Pachete utilitare -------------------------------------------------------
\usepackage{graphicx}
\usepackage{caption}
% Etichetele figurilor/tabelelor conform template-ului: Arial 10 pt, bold,
% centrate. Figura: titlu SUB figură, separator „–" („Figura 1 – ...").
% Tabelul: titlu DEASUPRA tabelului, separator „." („Tabelul 1. ...").
\captionsetup{font={footnotesize,bf},justification=centering}
\captionsetup[figure]{position=below,labelsep=endash}
\captionsetup[table]{position=above,labelsep=period}
% babel-romanian setează \tablename = „Tabela"; template-ul cere „Tabelul".
% Suprascriem la \begin{document} ca să batem definiția impusă de babel.
\AtBeginDocument{%
    \renewcommand{\figurename}{Figura}%
    \renewcommand{\tablename}{Tabelul}%
    \renewcommand{\listfigurename}{Lista figurilor}%
    \renewcommand{\listtablename}{Lista tabelelor}%
}

\usepackage{float}
\usepackage{array}
\usepackage{booktabs}
\usepackage{multirow}
\usepackage{longtable}

\usepackage{amsmath, amssymb, amsthm}
\usepackage{textcomp} % oferă \texteuro pentru simbolul €
\usepackage{siunitx}
\sisetup{output-decimal-marker={,}} % virgulă ca separator zecimal (RO)
% Unități non-standard folosite în lucrare — siunitx nu le definește implicit.
\DeclareSIUnit{\g}{g}          % accelerația gravitațională standard (context IMU)
\DeclareSIUnit{\byte}{B}       % octet (pentru rate de transfer)
\DeclareSIUnit{\LSB}{LSB}      % „least significant bit" (rezoluție senzor/ADC)
\DeclareSIUnit{\euro}{\text{\texteuro}}

\usepackage{listings}
\usepackage{xcolor}
\definecolor{codebg}{rgb}{0.97,0.97,0.97}
\definecolor{codecomment}{rgb}{0.4,0.4,0.4}
\definecolor{codekeyword}{rgb}{0.0,0.0,0.6}
\definecolor{codestring}{rgb}{0.6,0.0,0.0}
\lstset{
    backgroundcolor=\color{codebg},
    basicstyle=\ttfamily\footnotesize,
    breaklines=true,
    commentstyle=\color{codecomment}\itshape,
    keywordstyle=\color{codekeyword}\bfseries,
    stringstyle=\color{codestring},
    frame=single,
    framerule=0.3pt,
    rulecolor=\color{black!20},
    tabsize=2,
    showstringspaces=false,
    captionpos=b
}

% --- Comenzi personalizate (definite ÎNAINTE de hyperref, care le folosește
%     în \hypersetup pentru metadatele PDF) -------------------------------------
\newcommand{\candidatname}{Valentin-Gabriel NICOLA}        % <<< COMPLETAȚI
\newcommand{\coordinatorname}{Asist./Ș.L./Conf./Prof. dr. ing. Prenume NUME}
\newcommand{\sessionname}{Iunie 2026}            % <<< COMPLETAȚI
\newcommand{\specializare}{Automatica si Informatica Aplicata}
\newcommand{\workTitle}{Sistem portabil de cartografiere indoor cu LIDAR 2D și fuziune senzorială inerțială}

\usepackage{hyperref}
\hypersetup{
    colorlinks=true,
    linkcolor=black,
    citecolor=black,
    urlcolor=blue!50!black,
    pdftitle={\workTitle{}},
    pdfauthor={\candidatname{}}
}
\usepackage{url}

% --- TikZ pentru schemele derivate din cod ----------------------------------
\usepackage{tikz}
\usetikzlibrary{
    shapes.geometric,
    shapes.misc,
    arrows.meta,
    positioning,
    fit,
    calc,
    backgrounds,
    decorations.pathreplacing
}

% --- Bibliografie ------------------------------------------------------------
% Stil IEEE. Compilăm cu bibtex (mai compatibil cu Overleaf default).
\bibliographystyle{IEEEtran}

% Redenumim capitolul bibliografiei + îl adăugăm automat în Cuprins.
%   - \bibname e folosit ca titlu al capitolului produs de \bibliography{...}
%     iar formatarea Heading 1 (titlesec) îl va converti automat la majuscule.
%   - tocbibind pune intrarea în TOC fără să avem nevoie de un \chapter* manual
%     (evită capitolul "BIBLIOGRAFIE" dublu).
% Lista figurilor și lista tabelelor TREBUIE să apară în Cuprins (ca în
% lucrările-model), deci NU folosim opțiunile notlof/notlot.
\usepackage[nottoc,chapter]{tocbibind}
\renewcommand{\bibname}{Bibliografie}

% (Comenzile personalizate \candidatname / \coordinatorname / \sessionname /
%  \specializare / \workTitle sunt definite mai sus, înainte de hyperref.)

% Mediu pentru formule numerotate la dreapta cu paranteze rotunde
% (e.g. (1), (2), (3) — așa cum cere template-ul UPT)
% \begin{equation} ... \end{equation} face deja acest lucru implicit.

% Renunțăm la aspectul „monospace" pentru numele de pachete, fișiere și
% identificatori din text: arăta încărcat/„de cod" și diferit de lucrările-model.
% \texttt scrie de acum în fontul normal al textului (Arial), păstrând doar
% conținutul. \textnormal funcționează și în mod matematic (pentru numele de
% constante din formulele afișate).
\renewcommand{\texttt}[1]{\textnormal{#1}}

% =============================================================================
% ÎNCEPUTUL DOCUMENTULUI
% =============================================================================
\begin{document}

% Front matter: fără antet și fără număr de pagină vizibil (numerotarea curge
% totuși din pagina de titlu — vezi regula (e) din template). Comutăm la
% „fancy" abia la Introducere.
\pagestyle{coperta}

% --- Pagină de titlu --------------------------------------------------------
\thispagestyle{coperta}
\begin{titlepage}
    \centering

    % Sigla UPT, sus, centrată (ca în lucrările-model).
    {\includegraphics[height=2.2cm]{upt_logo.jpeg}\par}
    \vspace{0.6cm}

    % Identificarea instituției — scris normal, fără majuscule forțate.
    {\large\bfseries Universitatea Politehnica Timișoara\par}
    \vspace{0.15cm}
    Facultatea de Automatică și Calculatoare\par
    Domeniul: Calculatoare și Tehnologia Informației\par
    Specializarea: \specializare\par

    \vspace{2.5cm}

    {\itshape Lucrare de licență\par}
    \vspace{0.6cm}
    {\fontsize{18pt}{22pt}\bfseries \MakeUppercase{\workTitle{}}\par}

    \vfill

    \begin{flushleft}
        \textbf{Candidat:} \candidatname{}\par
        \vspace{0.3cm}
        \textbf{Coordonator științific:} \coordinatorname{}\par
    \end{flushleft}

    \vspace{1.2cm}

    Timișoara\par
    \sessionname{}\par
    \vspace{0.4cm}
\end{titlepage}

% =============================================================================
% REZUMAT
% =============================================================================
\thispagestyle{coperta}
% Titlu REZUMAT — Arial 20 pt, Bold, Uppercase, Center (per template).
\vspace*{12pt}
\begin{center}{\fontsize{20pt}{24pt}\bfseries REZUMAT}\end{center}
\addcontentsline{toc}{chapter}{Rezumat}
\vspace{24pt}

\noindent Lucrarea de față prezintă \textbf{Recon-Platform-R2}, un dispozitiv
portabil cu care se poate ridica planul 2D al unei încăperi. Ideea de
folosire este simplă: utilizatorul ține aparatul în mână și se plimbă prin
spațiu, iar harta pereților se construiește singură, în timp real, pe ecranul
unui telefon sau laptop. Aparatul are trei componente principale: un senzor
LIDAR 2D \textit{LD14P}, un calculator \textit{Raspberry Pi~5} și un
microcontroler \textit{ESP32} în rol de coprocesor.

Cartografierea propriu-zisă este făcută de algoritmul \textit{Karto SLAM}
(prin pachetul slam\_toolbox din \textit{ROS2 Jazzy Jalisco}), care produce o
hartă de tip \textit{occupancy grid}. Orientarea aparatului este estimată
dintr-un IMU \textit{MPU-6050}, prin integrarea vitezei unghiulare și un
filtru Kalman extins. Un element mai puțin obișnuit al sistemului este modul
în care ajung datele LIDAR la Pi: în loc de o conexiune serială directă, ele
trec prin ESP32, iar pe Pi sunt redate printr-un port serial „virtual" pe
care driverul de LIDAR îl deschide ca pe unul real.

Hărțile pot fi salvate într-o bază de date PostgreSQL și apoi curățate
printr-un set de operații de procesare a imaginii (filtru median, operații
morfologice, etichetarea componentelor conectate și transformata Hough),
care scot la iveală pereții ca segmente drepte. Toată interacțiunea cu
utilizatorul se face printr-o interfață web, accesibilă prin rețeaua WiFi
proprie a dispozitivului.

\medskip

\noindent\textbf{Cuvinte cheie:} SLAM, LIDAR 2D, ROS2 Jazzy Jalisco,
fuziune senzorială, EKF, ESP32, Raspberry Pi 5, transformata Hough,
PostgreSQL, Flask-SocketIO.

\clearpage

% =============================================================================
% ABSTRACT
% =============================================================================
\thispagestyle{coperta}
% Titlu ABSTRACT — Arial 20 pt, Bold, Uppercase, Center (per template).
\vspace*{12pt}
\begin{center}{\fontsize{20pt}{24pt}\bfseries ABSTRACT}\end{center}
\addcontentsline{toc}{chapter}{Abstract}
\vspace{24pt}

\selectlanguage{english}
\noindent This thesis presents \textbf{Recon-Platform-R2}, a handheld device
that produces a 2D floor plan of an indoor space. The idea of use is simple:
the operator holds the device and walks through the room, and the map of the
walls builds itself in real time on a phone or laptop screen. The device has
three main parts: a 2D \textit{LD14P} LIDAR sensor, a \textit{Raspberry Pi~5}
single-board computer, and an \textit{ESP32} microcontroller acting as a
coprocessor.

The mapping itself is done by the \textit{Karto SLAM} algorithm (through the
slam\_toolbox package on \textit{ROS2 Jazzy Jalisco}), which produces an
occupancy grid. The device orientation is estimated from an \textit{MPU-6050}
IMU, by integrating the angular rate and feeding it to an extended Kalman
filter. A less common feature of the system is how the LIDAR data reaches the
Pi: instead of a direct serial link, it travels through the ESP32 and is then
replayed on the Pi through a „virtual" serial port that the LIDAR driver opens
as if it were a real one.

Maps can be saved to a PostgreSQL database and later cleaned up with a set of
image-processing steps (median filter, morphological operations,
connected-component labelling, and the Hough transform) that turn the walls
into straight segments. All user interaction happens through a web interface,
reachable over the device's own WiFi access point.

\medskip

\noindent\textbf{Keywords:} SLAM, 2D LIDAR, ROS2 Jazzy Jalisco, sensor
fusion, EKF, ESP32, Raspberry Pi 5, Hough transform, PostgreSQL,
Flask-SocketIO.
\selectlanguage{romanian}

\clearpage

% =============================================================================
% CUPRINS
% =============================================================================
\thispagestyle{coperta}
\renewcommand{\contentsname}{CUPRINS}
\tableofcontents
\clearpage

% =============================================================================
% ELEMENTE SUPLIMENTARE — listă de abrevieri, lista figurilor, lista tabelelor
% (urmează structura lucrărilor de licență AC din anii anteriori)
% =============================================================================
\thispagestyle{coperta}
\vspace*{12pt}
\begin{center}{\fontsize{14pt}{17pt}\bfseries LISTĂ DE ABREVIERI}\end{center}
\addcontentsline{toc}{chapter}{Listă de abrevieri}
\vspace{12pt}

\begin{longtable}{@{}>{\bfseries}p{2.8cm} p{\dimexpr\textwidth-3.2cm\relax}@{}}
SLAM      & Simultaneous Localization and Mapping (localizare și cartografiere simultană) \\
LIDAR     & Light Detection and Ranging (detecție și telemetrie cu laser) \\
IMU       & Inertial Measurement Unit (unitate de măsurare inerțială) \\
EKF       & Extended Kalman Filter (filtru Kalman extins) \\
ROS2      & Robot Operating System 2 (middleware robotic) \\
DDS       & Data Distribution Service \\
QoS       & Quality of Service (calitatea serviciului) \\
SBC       & Single-Board Computer (calculator pe o singură placă) \\
UART      & Universal Asynchronous Receiver-Transmitter \\
USB-CDC   & USB Communications Device Class \\
PTY       & Pseudo-Terminal (pseudo-terminal Linux) \\
GPIO      & General-Purpose Input/Output \\
I\textsuperscript{2}C & Inter-Integrated Circuit (magistrală serială sincronă) \\
CRC       & Cyclic Redundancy Check (cod ciclic de verificare a integrității) \\
DLPF      & Digital Low-Pass Filter (filtru digital trece-jos) \\
CMOS      & Complementary Metal-Oxide-Semiconductor \\
NPN       & Tranzistor bipolar cu joncțiuni Negativ-Pozitiv-Negativ \\
PWM       & Pulse-Width Modulation (modulație în lățime de impuls) \\
LSB       & Least Significant Bit (bit cel mai puțin semnificativ) \\
TF2       & Sistemul de transformări de coordonate din ROS2 \\
ORM       & Object-Relational Mapping (cartografiere obiect-relațional) \\
REST      & Representational State Transfer \\
JSON      & JavaScript Object Notation \\
AP        & Access Point (punct de acces WiFi) \\
DHCP      & Dynamic Host Configuration Protocol \\
DNS       & Domain Name System \\
WPA2      & Wi-Fi Protected Access 2 \\
BFS       & Breadth-First Search (parcurgere în lățime) \\
LTS       & Long-Term Support (suport pe termen lung) \\
ARM       & Advanced RISC Machines (arhitectură de procesor) \\
SI        & Sistemul Internațional de unități de măsură \\
8N1       & Format serial: 8 biți de date, fără paritate, 1 bit de stop \\
\end{longtable}
\clearpage

% Lista figurilor și lista tabelelor (generate automat, adăugate în Cuprins).
\thispagestyle{coperta}
\listoffigures
\clearpage

\thispagestyle{coperta}
\listoftables
\clearpage

% =============================================================================
% CAPITOLUL 1 — INTRODUCERE
% =============================================================================
% De aici începe corpul lucrării: activăm antetul + numărul de pagină pe toate
% paginile. \let\ps@plain\ps@fancy forțează și paginile de început de capitol
% (care altfel ar folosi stilul „plain") să afișeze antetul complet.
\pagestyle{fancy}
\makeatletter
\let\ps@plain\ps@fancy
\makeatother

\chapter{Introducere}
\label{cap:introducere}

De multe ori avem nevoie de planul unei camere sau al unui apartament: când
ne mutăm, când vrem să rearanjăm mobila sau, pur și simplu, când un robot are
nevoie de o hartă a locului în care se mișcă. A obține acest plan nu este însă
mereu comod. Variantele existente sunt fie scumpe (scanerele laser
profesionale costă mii de euro), fie limitate (aplicațiile de pe telefon care
folosesc senzorul LiDAR au rază mică și depind de un anumit ecosistem).

Pornind de la această observație, lucrarea propune \textit{Recon-Platform-R2}:
un dispozitiv mic, ieftin și complet open-source, care produce harta 2D a unei
încăperi. Țelul nu este să concureze cu echipamentele profesionale, ci să arate
că, sub 200~EUR și folosind componente accesibile, se poate construi un scaner
portabil care funcționează singur, fără cabluri și fără rețea externă.

\section{Motivația alegerii temei}
\label{sec:motivatia}

% IDEI DE DEZVOLTAT:
%  - Carența soluțiilor open-source portabile (vs. proiectele de tip
%    rover autonom care necesită mecanică complexă)
%  - Pivotul de la robotul autonom „roomba" la dispozitivul portabil
%    (vezi commit-urile H1 și H1.5 — 2026-05-10)
%  - Utilizarea ROS2 ca middleware permite reutilizarea unei stive
%    mature de SLAM (slam_toolbox) și fuziune (robot_localization)
%    fără a reimplementa nimic.

Majoritatea proiectelor DIY de cartografiere bazate pe Raspberry~Pi și
senzori LIDAR ieftini pleacă de la un robot mobil: o platformă pe roți, cu
motoare, drivere de putere și o stivă de navigație. Mi s-a părut că aici e o
nepotrivire. Pentru a face o hartă nu am nevoie de roboți care se deplasează
singuri; roțile, encoderele și evitarea coliziunilor rezolvă problema
locomoției, nu pe cea a cartografierii. De ce să car după mine toată această
complexitate, când eu, ca om, mă pot deplasa prin cameră mult mai bine decât
orice rover?

De altfel, proiectul chiar a pornit ca un astfel de robot autonom (numit
intern \textit{roomba}): avea control de motoare, telecomandă printr-un
\textit{gamepad} Bluetooth și un modul de explorare automată. Pe măsură ce
lucram la el, am observat că aproape tot efortul se ducea în partea mecanică
— motoare, alimentare, fiabilitatea tracțiunii — în timp ce partea care mă
interesa cu adevărat (SLAM, fuziunea senzorilor, vizualizarea hărții) nu avea
nicio legătură cu modul în care se mișcă aparatul. Așa că am renunțat la
locomoție și am transformat robotul într-un dispozitiv pe care îl porți în
mână, păstrând doar senzorii și partea de calcul.

Pentru partea software am ales \textit{ROS2} ca platformă de bază. Avantajul
este că pot reutiliza componente mature, deja testate de comunitate — SLAM
(slam\_toolbox) și fuziune senzorială (robot\_localization) — fără să
reimplementez algoritmi consacrați. Astfel m-am putut concentra pe ceea ce
era cu adevărat specific proiectului: integrarea pieselor, modul de transport
al datelor de la LIDAR și interfața cu utilizatorul.

\section{Obiectivele lucrării}
\label{sec:obiective}

% IDEI DE DEZVOLTAT:
%  - Obiectiv principal: hartă 2D recognoscibilă a unei camere de
%    5x5m dintr-o singură plimbare
%  - Obiective specifice numerice (acceptance criteria din ROADMAP.md):
%      * Erori sub ±10 cm pe pereții lungi
%      * Fără avertismente "scan match failed" la mers normal
%      * Salvare pe disc + reîncărcare prin /api/maps/<id>/data
%  - Obiective non-funcționale: cost <200 EUR, autonomie minim 1 oră,
%    operabil de o persoană fără instruire tehnică

Obiectivul principal este destul de concret: vreau ca, dintr-o singură
plimbare prin cameră, dispozitivul să producă harta 2D a pereților unei
încăperi de aproximativ $5 \times \SI{5}{\meter}$. Pentru a considera acest
obiectiv atins, mi-am stabilit câteva criterii măsurabile:

\begin{itemize}\itemsep0pt
    \item închidere de buclă vizibilă pe traseul perimetral;
    \item erori dimensionale sub $\pm\SI{10}{\centi\meter}$ pe pereții
          lungi (validat prin măsurarea cu ruleta);
    \item lipsa avertismentelor \texttt{scan match failed} în jurnalele
          \texttt{slam\_toolbox} la o cadență de mers de
          $\sim\SI{0.5}{\meter\per\second}$;
    \item harta salvată în baza de date trebuie să poată fi reîncărcată
          identic prin endpoint-ul REST \texttt{/api/maps/<id>/data}.
\end{itemize}

Pe lângă acestea, există și câteva cerințe care țin de utilizare, nu de
acuratețe: costul total al componentelor să rămână sub $\SI{200}{\euro}$,
aparatul să funcționeze cel puțin o oră pe baterie și să poată fi folosit de
oricine, fără pregătire tehnică — în mod ideal, apăsând un singur buton
\textit{Start scan} în pagina web.

\section{Structura lucrării}
\label{sec:structura}

% IDEI DE DEZVOLTAT:
%  - Maparea capitole -> componente proiect:
%    cap.2 = teorie (algoritmi); cap.3 = hardware; cap.4 = protocol UART;
%    cap.5 = firmware; cap.6 = software Pi; cap.7 = web; cap.8 = DB;
%    cap.9 = integrare; cap.10 = testare; cap.11 = concluzii.

Lucrarea este structurată în 12 capitole. \textbf{Capitolul~2} oferă o
trecere în revistă a stadiului actual al cartografierii LIDAR portabile
și a soluțiilor concurente. \textbf{Capitolul~3} prezintă fundamentele
teoretice ale algoritmilor folosiți: SLAM bazat pe optimizarea grafului
de poze, filtre Kalman extinse, transformata Hough probabilistică,
operații morfologice, etichetarea componentelor conectate, integrarea
giroscopică și verificarea integrității cu CRC-8. \textbf{Capitolele~4--10}
descriu implementarea sistemului în straturi succesive: arhitectura
hardware, protocolul de comunicație ESP32~$\leftrightarrow$~Pi,
firmware-ul ESP32, stiva software ROS2 de pe Pi, interfața web,
persistența în baza de date și integrarea finală.
\textbf{Capitolul~11} prezintă strategia de testare și rezultatele
validării. \textbf{Capitolul~12} sintetizează contribuțiile și
direcțiile viitoare.

\clearpage

% =============================================================================
% CAPITOLUL 2 — STADIUL ACTUAL
% =============================================================================
\chapter{Stadiul actual al tehnologiei}
\label{cap:stadiul}

% IDEI DE DEZVOLTAT:
%   Comparație cu:
%   - Scanere comerciale (Leica BLK2GO, Kaarta Stencil) — preț, acuratețe
%   - Aplicații smartphone (iPhone Pro LiDAR scanner, Polycam) — limitări
%   - Proiecte DIY open-source (rPLIDAR + RPi cu Hector SLAM)
%   - Roboti autonomi (TurtleBot, Pioneer) ca alternativă

\section{Sisteme SLAM 2D existente}
\label{sec:slamexistente}

Problema localizării și cartografierii simultane (SLAM) este una dintre
cele mai studiate din robotica mobilă~\cite{thrun2005probrobotics,
lidar_slam_survey}. În regimul bidimensional, bazat pe scanări LIDAR, s-au
maturizat mai multe familii de algoritmi, disponibili ca pachete
open-source integrate în ROS/ROS2:

\begin{itemize}\itemsep0pt
    \item \textbf{GMapping} --- abordare bazată pe filtre particulare
          Rao-Blackwellizate. Robustă pe trasee scurte, dar costisitoare
          ca memorie (un set de hărți per particulă) și fără închidere de
          buclă globală.
    \item \textbf{Hector SLAM}~\cite{kohlbrecher2011hector} --- potrivire
          de scanări pe o hartă multi-rezoluție, fără odometrie. Excelează
          la frecvențe LIDAR mari, dar derapează în medii cu puține
          trăsături geometrice deoarece nu menține un graf de poze.
    \item \textbf{Cartographer}~\cite{hess2016cartographer} (Google) ---
          SLAM bazat pe sub-hărți și optimizare a grafului de poze, cu
          închidere de buclă prin potrivire \textit{branch-and-bound}.
          Foarte precis, dar greu de acordat și mai solicitant
          computațional.
    \item \textbf{Karto SLAM}~\cite{konolige2010karto} --- SLAM bazat pe
          optimizarea grafului de poze cu un rezolvitor sparse de cele mai
          mici pătrate. Este nucleul algoritmic al pachetului
          \texttt{slam\_toolbox}~\cite{macenski2021slamtoolbox}.
\end{itemize}

Pachetul \texttt{slam\_toolbox} a fost ales pentru acest proiect deoarece
este menținerea de referință pentru SLAM 2D în ROS2, oferă un mod
\textit{online asincron} potrivit pentru rularea în timp real pe un
sistem cu resurse limitate, dispune de închidere de buclă și de
serializarea/deserializarea grafului de poze, și expune parametri de
acordare bogați (Tabelul~\ref{tab:slam-compare}). Tot el oferă servicii
ROS (pauză, resetare) folosite direct de interfața web a proiectului.

\begin{table}[H]
\centering
\caption{Comparație sintetică a algoritmilor SLAM 2D open-source relevanți.}
\label{tab:slam-compare}
\renewcommand{\arraystretch}{1.15}
\begin{tabular}{lcccc}
\toprule
\textbf{Algoritm} & \textbf{Tip} & \textbf{Închidere buclă} & \textbf{Cost CPU} & \textbf{ROS2} \\
\midrule
GMapping       & Filtru particular  & Nu       & Mediu  & Parțial \\
Hector SLAM    & Scan-matching      & Nu       & Scăzut & Parțial \\
Cartographer   & Sub-hărți + graf   & Da       & Ridicat& Da \\
Karto / \texttt{slam\_toolbox} & Graf de poze & Da & Mediu  & Da (nativ) \\
\bottomrule
\end{tabular}
\end{table}

\section{Senzori LIDAR low-cost pentru cartografierea indoor}
\label{sec:lidarexistente}

Piața senzorilor LIDAR rotativi low-cost destinați aplicațiilor hobby și
educaționale s-a diversificat considerabil în ultimii ani. Cei mai
răspândiți candidați pentru cartografierea indoor sunt:

\begin{itemize}\itemsep0pt
    \item \textbf{SLAMTec RPLIDAR A1/A2/A3} --- 360°, rază
          $\SI{6}{}$--$\SI{25}{\meter}$ în funcție de model, interfață
          UART/USB; A1 este cel mai ieftin ($\sim\SI{100}{\euro}$) dar are
          rezoluție unghiulară modestă.
    \item \textbf{YDLIDAR X4/X2} --- 360°, rază $\sim\SI{10}{\meter}$,
          frecvență de scanare reglabilă, preț comparabil cu RPLIDAR~A1.
    \item \textbf{LDROBOT LD14P (LD-D200)}~\cite{ld14p} --- 360°, rază
          $\SI{0.02}{}$--$\SI{8}{\meter}$, frecvență nominală
          $\SI{10}{\hertz}$, interfață UART la $\SI{230400}{}$ baud,
          ieșire CMOS la $\SI{3.3}{\volt}$, cost $\sim\SI{50}{\euro}$.
\end{itemize}

Senzorul \textbf{LD14P} a fost ales pentru raportul preț/performanță cel
mai favorabil din segment și pentru interfața UART simplă, la nivel de
octet, cu ieșire logică direct la $\SI{3.3}{\volt}$ --- compatibilă cu
intrările ESP32 fără translator de nivel. Protocolul său la nivel de cadru
este comun cu seria LD19 (motiv pentru care driverul
\texttt{ldlidar\_stl\_ros2} este lansat cu parametrul
\texttt{product\_name: LDLiDAR\_LD19}). Limitarea sa principală --- raza
de doar $\SI{8}{\meter}$ --- este acceptabilă pentru spațiile interioare
domestice țintă (camere, apartamente, birouri mici). Particularitatea
electrică a senzorului (firul \textit{verde} de RX legat la masa locală
prin tranzistor pentru selecția vitezei implicite, vezi
§\ref{sec:lidar-wires}) a fost esențială pentru proiectarea circuitului
de control al motorului.

\section{Fuziune senzorială inerțială pentru estimarea pozițiilor}
\label{sec:imuexistente}

Estimarea orientării dintr-un IMU cu 6 axe (accelerometru + giroscop, fără
magnetometru) se face în mod uzual cu un filtru complementar care fuzează
informația de gravitație din accelerometru (corectează \textit{roll} și
\textit{pitch}) cu integrarea vitezei unghiulare din giroscop. Cele mai
cunoscute astfel de filtre sunt:

\begin{itemize}\itemsep0pt
    \item \textbf{Filtrul Madgwick}~\cite{madgwick2014} --- bazat pe
          descreșterea gradientului unei funcții de eroare a orientării,
          cu un singur parametru de acordare (\textit{gain});
    \item \textbf{Filtrul Mahony}~\cite{mahony2008} --- filtru
          complementar neliniar pe grupul $SO(3)$, cu termeni proporțional
          și integral;
    \item \textbf{Filtrul Kalman extins generalizat}~\cite{moore2014robotlocalization}
          --- abordarea adoptată de pachetul \texttt{robot\_localization}.
\end{itemize}

În proiect am ales totuși o cale mai simplă: integrez direct viteza unghiulară
$\omega_z$ a giroscopului, fără să corectez cu accelerometrul. Motivul ține de
o particularitate a cipului MPU-6050 folosit. Acesta are un bias de fabrică
(\texttt{ZA\_OFFSET}) care face ca, în repaus, axa Z a accelerometrului să
arate aproximativ $\SI{15.3}{\meter\per\second\squared}$ în loc de
$\SI{9.81}{}$~\cite{mpu6050regmap}. Cu alte cuvinte, vectorul de gravitație
măsurat are direcția corectă, dar mărimea greșită, așa că un filtru de tip
Madgwick sau Mahony — care se bazează tocmai pe „gravitația arată în jos" — ar
introduce o eroare constantă în înclinare.

Din fericire, pentru un scaner ținut în mână, înclinarea contează prea puțin:
scanarea se face practic într-un plan orizontal, deci singurul unghi care mă
interesează este cel de cap (\textit{yaw}), care oricum se obține numai din
giroscop (cipul nu are magnetometru). Giroscopul are o derivă lentă (în jur de
$\SI{0.7}{\degree\per\second}$ pe exemplarul testat), însă aceasta este
corectată permanent de potrivirea de scanări din slam\_toolbox. De aceea, o
integrare simplă este suficientă în practică; trecerea la un filtru Madgwick,
cu o calibrare prealabilă a biasului, rămâne ca îmbunătățire viitoare.

\section{Soluții comerciale concurente}
\label{sec:comerciale}

La capătul opus al pieței se află scanerele mobile profesionale, care
combină LIDAR 3D, IMU de grad industrial și, uneori, camere pentru
colorarea norului de puncte. Reprezentative sunt \textbf{Leica BLK2GO},
\textbf{Kaarta Stencil}, \textbf{GeoSLAM ZEB Horizon} și \textbf{Faro
Focus}. Aceste sisteme produc nori de puncte 3D de înaltă densitate, cu
acuratețe de ordinul centimetrilor sau milimetrilor, însă la prețuri
cuprinse între câteva mii și câteva zeci de mii de euro --- de două-trei
ordine de mărime peste bugetul vizat de proiectul de față.

O a doua categorie de concurenți o reprezintă aplicațiile pentru
\textbf{smartphone} care exploatează senzorul LiDAR ToF integrat în
dispozitivele Apple din gama Pro (de exemplu \textit{Polycam},
\textit{Scaniverse}). Acestea sunt accesibile, dar au rază utilă mică,
acuratețe limitată la distanță și sunt legate de ecosistemul comercial
iOS/iPadOS.

\begin{table}[H]
\centering
\caption{Poziționarea Recon-Platform-R2 față de soluțiile concurente.}
\label{tab:competitors}
\renewcommand{\arraystretch}{1.15}
\begin{tabular}{lccc}
\toprule
\textbf{Soluție} & \textbf{Ieșire} & \textbf{Acuratețe} & \textbf{Cost orientativ} \\
\midrule
Leica BLK2GO     & Nor 3D & mm--cm & $> \SI{50000}{\euro}$ \\
GeoSLAM ZEB Horizon & Nor 3D & cm   & $\sim\SI{30000}{\euro}$ \\
iPhone Pro + Polycam & Mesh 3D & dm  & cost telefon \\
\textbf{Recon-Platform-R2} & Hartă 2D & $\sim\SI{10}{\centi\meter}$ & $< \SI{200}{\euro}$ \\
\bottomrule
\end{tabular}
\end{table}

Concluzia comparației este că \textit{Recon-Platform-R2 nu concurează cu
produsele profesionale}, ci ocupă în mod deliberat nișa
,,scanare 2D suficient de bună la cost foarte redus și complet
open-source'', acolo unde un plan de etaj aproximativ este mai valoros
decât un nor de puncte 3D scump.

\clearpage

% =============================================================================
% CAPITOLUL 3 — FUNDAMENTE TEORETICE
% =============================================================================
\chapter{Fundamente teoretice}
\label{cap:teorie}

% IDEI DE DEZVOLTAT:
%   Acest capitol este "ancora teoretică" — fiecare algoritm folosit
%   în implementare are aici 2-4 pagini de fundament matematic.
%   Toate referințele bibliografice principale apar prima dată aici.

\section{Algoritmul Karto SLAM și optimizarea grafului de poze}
\label{sec:karto}

\textit{slam\_toolbox}~\cite{macenski2021slamtoolbox} folosește algoritmul
\textbf{Karto SLAM}~\cite{konolige2010karto}, un sistem 2D bazat pe
optimizarea unui graf de poze. Diagrama logică a algoritmului este:

\begin{enumerate}\itemsep0pt
    \item achiziție scanare LIDAR (\textit{cadru curent});
    \item potrivire (\textit{scan matching}) a cadrului curent față de
          ultimele $k$ cadre cheie (\textit{keyframes});
    \item dacă mișcarea de la ultimul cadru cheie depășește un prag, se
          adaugă cadrul curent ca nou nod în graf;
    \item căutare de potriviri pentru închidere de buclă (\textit{loop
          closure}) cu cadre cheie distante geometric;
    \item optimizare a întregului graf cu un rezolvitor de probleme
          neliniare ale celor mai mici pătrate
          (\textit{Ceres}~\cite{ceres_solver}).
\end{enumerate}

Ce înseamnă, mai exact, ,,localizare și cartografiere simultană''? Pe scurt,
dispozitivul trebuie să afle în același timp două lucruri care depind unul de
celălalt: pe unde a trecut (traiectoria $\mathbf{x}_{1:t}$) și cum arată harta
$\mathbf{m}$. Dacă aș ști harta, aș putea spune ușor unde mă aflu; dacă aș ști
unde mă aflu, aș putea construi ușor harta. Problema e că nu le știu pe niciuna
și trebuie să le estimez pe amândouă deodată, pornind de la observațiile
$\mathbf{z}_{1:t}$ și de la odometria $\mathbf{u}_{1:t}$~\cite{thrun2005probrobotics}:

\begin{equation}\label{eq:slam-posterior}
    p(\mathbf{x}_{1:t}, \mathbf{m} \mid \mathbf{z}_{1:t}, \mathbf{u}_{1:t}).
\end{equation}

Karto rezolvă această problemă transformând-o într-una de optimizare, iar
ideea din spate este destul de intuitivă. Fiecare poziție din care s-a făcut o
scanare devine un \textit{nod} într-un graf. Între noduri se trag \textit{muchii}
care spun, în esență, ,,din poziția $i$ până în poziția $j$ aparatul s-a mutat
cu atât'' — informație obținută prin suprapunerea celor două scanări. Fiecare
muchie are atașată și o ,,încredere'' (matricea de informație
$\boldsymbol{\Omega}_{ij}$): cu cât potrivirea a fost mai clară, cu atât
constrângerea cântărește mai mult. Soluția căutată este setul de poziții care
se împacă cel mai bine cu toate aceste constrângeri, adică minimizează suma
erorilor (ecuația~\eqref{eq:slam-cost}). Dacă presupunem că zgomotul de
măsurare este gaussian, acest cel-mai-mic-pătrat coincide chiar cu cea mai
probabilă traiectorie și hartă din ecuația~\eqref{eq:slam-posterior}.

Eroarea unei muchii (ecuația~\eqref{eq:slam-residual}) măsoară cât de mult
diferă deplasarea măsurată $\mathbf{z}_{ij}$ de cea prezisă de pozițiile
curente:

\begin{equation}\label{eq:slam-residual}
    \mathbf{e}_{ij}(\mathbf{x}_i, \mathbf{x}_j, \mathbf{z}_{ij})
        = \mathbf{z}_{ij} \ominus \big(\mathbf{x}_i^{-1} \oplus \mathbf{x}_j\big),
\end{equation}

\noindent unde $\oplus$ și $\ominus$ sunt compunerea și, respectiv, diferența
a două transformări rigide din plan ($SE(2)$). Sistemul neliniar care rezultă
este rezolvat iterativ de către \textit{Ceres}~\cite{ceres_solver}, configurat
în \texttt{slam\_params.yaml}. Marele câștig al acestei formulări este
\textit{închiderea de buclă}: când aparatul revine într-un loc pe unde a mai
trecut, se adaugă o muchie nouă între două cadre îndepărtate în timp, dar
asemănătoare ca aspect; optimizarea ,,trage'' apoi întreaga traiectorie înapoi
pe poziția corectă, anulând deriva acumulată între timp.

% Formula 1 — costul SLAM ca optimizare a grafului de poze
\begin{equation}\label{eq:slam-cost}
    \mathbf{x}^{*} = \arg\min_{\mathbf{x}}
        \sum_{(i,j) \in \mathcal{E}}
            \left\| \mathbf{e}_{ij}(\mathbf{x}_i, \mathbf{x}_j, \mathbf{z}_{ij}) \right\|^{2}_{\boldsymbol{\Omega}_{ij}}
\end{equation}

\noindent unde $\mathbf{x}_i$ sunt pozele estimate, $\mathbf{z}_{ij}$
sunt măsurătorile relative dintre cadrele cheie $i$ și $j$,
$\boldsymbol{\Omega}_{ij}$ este matricea informației (inversa
covarianței), iar $\mathbf{e}_{ij}$ este eroarea reziduală.

\section{Filtre Kalman extinse pentru fuziune senzorială}
\label{sec:ekf}

Filtrul Kalman~\cite{kalman1960} este, pe scurt, o metodă de a combina o
predicție cu o măsurătoare nouă și de a obține o estimare mai bună decât
fiecare luată separat. Funcționează optim pentru sisteme liniare cu zgomot
gaussian; varianta \textit{extinsă} (EKF) îl duce și la sisteme neliniare,
liniarizând local funcțiile de tranziție $\mathbf{f}$ și de observație
$\mathbf{h}$ în jurul estimării curente~\cite{welchbishop2006kalman}. Filtrul
lucrează în doi pași care se repetă la fiecare măsurătoare:
\textit{predicția}, care duce starea înainte în timp conform modelului
(ecuația~\eqref{eq:ekf-predict}), și \textit{actualizarea}, care o corectează
cu noua măsurătoare $\mathbf{z}_k$, cântărind cele două surse prin câștigul
Kalman $\mathbf{K}_k$ (ecuațiile~\eqref{eq:ekf-update} și~\eqref{eq:kalman-gain}).
Incertitudinea estimării se propagă prin jacobianul
$\mathbf{F}_k = \partial\mathbf{f}/\partial\mathbf{x}$:

\begin{equation}\label{eq:ekf-cov}
    \mathbf{P}_{k|k-1} = \mathbf{F}_k\,\mathbf{P}_{k-1|k-1}\,\mathbf{F}_k^T + \mathbf{Q}_k,
\end{equation}

\noindent unde $\mathbf{Q}_k$ este covarianța zgomotului de proces, iar
$\mathbf{R}_k$ (din ecuația~\eqref{eq:kalman-gain}) covarianța zgomotului
de măsurare.

În proiect folosesc implementarea de EKF din pachetul
robot\_localization~\cite{moore2014robotlocalization}, configurată
(în \texttt{config/ekf.yaml}) în mod 2D — adică ignoră tot ce iese din planul
orizontal. Mai mult, las filtrul să folosească din IMU \textbf{doar} unghiul
de cap $\theta$ (\textit{yaw}) și viteza lui de variație $\dot{\theta}$;
celelalte 13 mărimi posibile sunt oprite. Practic, starea filtrului se reduce
la

\begin{equation}\label{eq:ekf-state}
    \mathbf{x} = \begin{bmatrix} \theta & \dot{\theta} \end{bmatrix}^{T}.
\end{equation}

Accelerația liniară o las complet pe dinafară, din motivul explicat în
§\ref{sec:imuexistente}: biasul \texttt{ZA\_OFFSET} ar fi integrat de două ori
și ar duce la o derivă de poziție scăpată de sub control. Așadar, poziția
aparatului nu vine din IMU, ci din corecția pe care o oferă slam\_toolbox la
fiecare potrivire de scanare. EKF-ul are aici un rol modest, dar util: să dea
o orientare netedă, care servește ca punct de pornire pentru potrivirea
scanărilor.

% Formula 2 — pas predicție EKF
\begin{equation}\label{eq:ekf-predict}
    \hat{\mathbf{x}}_{k|k-1} = \mathbf{f}(\hat{\mathbf{x}}_{k-1|k-1}, \mathbf{u}_k)
\end{equation}

% Formula 3 — pas actualizare EKF
\begin{equation}\label{eq:ekf-update}
    \hat{\mathbf{x}}_{k|k} = \hat{\mathbf{x}}_{k|k-1}
        + \mathbf{K}_k \left(\mathbf{z}_k - \mathbf{h}(\hat{\mathbf{x}}_{k|k-1})\right)
\end{equation}

\noindent cu câștigul Kalman:

\begin{equation}\label{eq:kalman-gain}
    \mathbf{K}_k = \mathbf{P}_{k|k-1} \mathbf{H}_k^T
        \left(\mathbf{H}_k \mathbf{P}_{k|k-1} \mathbf{H}_k^T + \mathbf{R}_k\right)^{-1}
\end{equation}

\section{Integrarea giroscopică și conversia în cuaternioni}
\label{sec:giro}

Înainte ca EKF-ul să primească o orientare, aceasta trebuie obținută din
datele brute ale giroscopului. Pasul este foarte simplu: pentru a afla cu cât
s-a rotit aparatul, însumez în timp viteza unghiulară. Cum scanarea se face
într-un plan orizontal, mă interesează numai rotația în jurul axei verticale,
deci doar componenta $\omega_z$:

\begin{equation}\label{eq:yaw-integrate}
    \theta_n = \theta_{n-1} + \omega_{z,n} \cdot \Delta t
\end{equation}

slam\_toolbox nu primește însă unghiul direct, ci sub formă de cuaternion
(așa cere mesajul standard \texttt{sensor\_msgs/Imu}); pentru o rotație doar
în plan, conversia este imediată:

\begin{equation}\label{eq:yaw-to-quat}
    \mathbf{q} = \left( 0,\; 0,\; \sin\left(\frac{\theta}{2}\right),\; \cos\left(\frac{\theta}{2}\right) \right)
\end{equation}

În cod, această însumare are câteva precauții practice. Pasul de timp
$\Delta t$ se calculează din diferența marcajelor de timp ale mesajelor; dacă
această diferență iese aberantă (negativă sau mai mare de o secundă, semn că
s-au pierdut pachete), pasul respectiv este ignorat, ca să nu apară salturi.
La final, unghiul este readus în intervalul $(-\pi, \pi]$.

Punctul slab al metodei este, evident, \textbf{deriva}: orice mic bias al
giroscopului se adună în timp, iar unghiul ,,fuge'' încet (am măsurat în jur de
$\SI{0.7}{\degree\per\second}$ pe exemplarul de test). În alt context, atât ar
fi o problemă; aici nu, pentru că slam\_toolbox recalculează orientarea din
suprapunerea scanărilor și, practic, șterge deriva acumulată între două
potriviri. Avem deci o împărțire elegantă a muncii: giroscopul dă o orientare
netedă și rapidă (la \SI{100}{\hertz}), iar SLAM-ul o readuce periodic la
realitate (la frecvența scanărilor, $\sim\SI{6}{}$--$\SI{10}{\hertz}$). Tocmai
de aceea o simplă integrare, fără filtru sofisticat, este suficientă în
practică.

\section{Transformata Hough probabilistică}
\label{sec:hough}

Etapa a 5-a a pipeline-ului de procesare a hărților salvate
(\texttt{recon\_db.\allowbreak postprocess.\allowbreak hough\_line\_segments}) folosește o
variantă probabilistică a transformatei Hough~\cite{hough1962, dudahart1972,
matas2000probhough} pentru a detecta segmente de pereți în harta curățată.

Parametrizarea polară a unei linii:

\begin{equation}\label{eq:hough-rho}
    \rho = x \cos\theta + y \sin\theta
\end{equation}

\noindent unde perechea $(\rho, \theta)$ identifică în mod unic o dreaptă.
Ideea transformatei Hough este surprinzător de simplă: fiecare pixel ocupat
,,votează'' pentru toate dreptele care ar putea trece prin el. Dreptele reale,
pe care stau mulți pixeli (de exemplu un perete), adună multe voturi și ies în
evidență; restul rămân cu voturi puține. Varianta \textit{probabilistică}
merge un pas mai departe: în loc de drepte infinite, urmărește dreapta votată
și păstrează din ea doar bucățile acoperite efectiv de pixeli — segmentele de
perete —, tolerând mici întreruperi.

În implementarea mea, unghiul $\theta$ este împărțit în 180 de pași (un grad
fiecare), se aleg direcțiile cu cele mai multe voturi, iar de-a lungul lor se
decupează segmentele reale de perete. Acestea sunt apoi desenate în interfața
web peste harta curățată, ceea ce face ca rezultatul să semene cu un plan de
etaj, nu cu o pată colorată aleatoriu.

\section{Operațiuni morfologice pe imagini binare}
\label{sec:morfologie}

Etapele 2 și 3 ale pipeline-ului aplică operațiuni morfologice clasice
(\textit{eroziune}, \textit{dilatare}, \textit{deschidere},
\textit{închidere})~\cite{serra1982morpho, soille2003morpho} pe masca
binară de celule ocupate.

\begin{align}
    \text{Dilatare:}\quad A \oplus B &= \{ z \mid (\hat{B})_z \cap A \neq \emptyset \} \label{eq:dilation}\\
    \text{Eroziune:}\quad A \ominus B &= \{ z \mid (B)_z \subseteq A \} \label{eq:erosion}
\end{align}

Intuitiv, \textbf{dilatarea} ,,îngroașă'' zonele ocupate, iar
\textbf{eroziunea} le ,,subțiază''. Combinându-le, obținem două operații utile:
\textbf{deschiderea} (mai întâi eroziune, apoi dilatare) șterge punctele
izolate de zgomot, iar \textbf{închiderea} (mai întâi dilatare, apoi eroziune)
astupă găurile mici dintr-un perete. În implementarea pur-NumPy, acestea se
obțin simplu: se suprapune masca binară cu cele opt deplasări vecine și se
combină prin SAU logic (pentru dilatare) sau prin ȘI logic (pentru eroziune).

Aici am dat însă peste o capcană, prinsă de altfel și de testele automate: o
deschidere aplicată chiar și o singură dată \textbf{șterge complet} pereții
mai subțiri de trei celule. Cum la rezoluția obișnuită (\SI{5}{\centi\meter}
pe celulă) un perete are deseori doar una-două celule grosime, deschiderea
l-ar face să dispară. De aceea, în mod implicit, las activă doar închiderea, iar
zgomotul punctual îl elimin în etapa următoare, prin filtrul de mărime minimă a
componentelor conectate (§\ref{sec:cc}) — care scoate o celulă rătăcită fără
să atingă pereții.

\section{Etichetarea componentelor conectate}
\label{sec:cc}

Următorul pas grupează celulele ocupate în ,,obiecte'' separate: două celule
ocupate care se ating (inclusiv pe diagonală) fac parte din aceeași
componentă~\cite{rosenfeldpfaltz1966}. Fiecare componentă primește o etichetă,
iar cele prea mici — sub \texttt{min\_cluster\_size} celule (implicit 8) —
sunt considerate zgomot și aruncate.

Pentru a găsi componentele, parcurg harta și, de fiecare dată când dau peste o
celulă ocupată neetichetată, pornesc o ,,umplere prin inundare'' care marchează
toate celulele lipite de ea înainte de a merge mai departe. Am scris această
parcurgere cu o stivă proprie, nu recursiv. Motivul e practic: o variantă
recursivă ar depăși ușor limita de recursie din Python pe o componentă mare
(un perete lung al camerei) și ar arunca o eroare. Cu stiva proprie nu mai
există această limită. La final, celulele rămase fără o componentă suficient de
mare sunt readuse la ,,liber'' — așa scap de artefactele potrivirii de scanări
și de obstacolele trecătoare, de pildă o persoană care a traversat o singură
dată câmpul LIDAR-ului.

\section{Filtrul median 3\texorpdfstring{$\times$}{x}3}
\label{sec:median}

Etapa~1 aplică (opțional) un filtru median 3$\times$3~\cite{tukey1977eda}
pentru a elimina celulele rătăcite. Operația este definită ca:

\begin{equation}\label{eq:median}
    g(x, y) = \operatorname{median}\{f(x+i, y+j) \mid (i, j) \in W\}
\end{equation}

\noindent unde $W = \{-1, 0, +1\}^2$ este fereastra 3$\times$3.

\section{Verificarea integrității datelor cu CRC-8 Dallas/Maxim}
\label{sec:crc}

Pe o legătură serială se mai întâmplă ca un octet să fie corupt pe drum. Ca să
prind astfel de erori, fiecare cadru trimis între ESP32 și Raspberry~Pi poartă
la final o sumă de control CRC-8, calculată cu polinomul
Dallas/Maxim~\cite{maxim1wirecrc, linuxcrc8}:

\begin{equation}\label{eq:crc-poly}
    p(x) = x^8 + x^5 + x^4 + 1
\end{equation}

\noindent scris pe scurt ca $\mathtt{0x07}$. La recepție, fiecare parte
recalculează aceeași sumă peste octeții primiți și o compară cu cea din cadru;
dacă nu coincid, cadrul este pur și simplu aruncat. Am ales acest polinom
pentru că este același cu cel folosit de utilitarul \texttt{crc8} din nucleul
Linux — astfel am putut verifica ușor că cele două implementări, cea de pe
ESP32 și cea de pe Pi, dau exact aceleași valori.

\clearpage

% =============================================================================
% CAPITOLUL 4 — ARHITECTURA HARDWARE
% =============================================================================
\chapter{Arhitectura hardware a sistemului}
\label{cap:hardware}

\section{Privire de ansamblu}
\label{sec:hw-overview}

% Schema 1 — arhitectura sistemului
\begin{figure}[H]
\centering
\begin{tikzpicture}[
    node distance=0.7cm,
    every node/.style={font=\footnotesize},
    block/.style={rectangle, draw, rounded corners, minimum width=2.5cm, minimum height=0.9cm, align=center, fill=blue!5},
    sensor/.style={rectangle, draw, minimum width=2cm, minimum height=0.8cm, align=center, fill=orange!10},
    arr/.style={-{Latex[length=2mm]}, thick}
]
    \node[sensor] (lidar) {LD14P\\LIDAR};
    \node[sensor, right=1.2cm of lidar] (imu) {MPU-6050\\IMU};
    \node[sensor, below=of imu] (btn) {3 butoane};
    \node[block, below=1.3cm of lidar] (esp) {ESP32 DevKit V1};
    \node[block, below=of esp, minimum width=6cm, minimum height=1.5cm] (pi) {Raspberry Pi 5\\ROS2 Jazzy + Flask};
    \node[block, right=1cm of pi, fill=green!10] (db) {PostgreSQL\\Docker};
    \node[block, below=of pi, fill=red!10] (browser) {Browser\\http://recon.local};

    \draw[arr] (lidar) -- node[left]{\scriptsize UART 230400} (esp);
    \draw[arr] (imu) -- node[right]{\scriptsize I\textsuperscript{2}C 400 kHz} (esp);
    \draw[arr] (btn) -- (esp);
    \draw[arr] (esp) -- node[right]{\scriptsize USB-CDC 460800} (pi);
    \draw[arr] (pi) -- (db);
    \draw[arr] (pi) -- node[right]{\scriptsize WiFi AP \texttt{Recon}} (browser);
\end{tikzpicture}
\caption{Arhitectura hardware de nivel înalt. Toate datele senzoriale (LIDAR, IMU, butoane) tranzitează ESP32-ul înainte de a ajunge la Raspberry~Pi prin USB-CDC.}
\label{fig:hw-overview}
\end{figure}

Ideea principală a montajului, vizibilă în Figura~\ref{fig:hw-overview}, este
că \textbf{toate} datele de la senzori intră în Raspberry~Pi printr-o
\textbf{singură} legătură USB. Inclusiv scanările LIDAR, care în mod normal
s-ar lega direct la portul serial al Pi-ului, trec acum prin ESP32 și ajung pe
aceeași cale ca eșantioanele de la IMU și apăsările de buton. Avantajul este
că Pi-ul are de gestionat un singur port, fără cabluri și interfețe în plus și
fără bătăile de cap legate de permisiunile pe portul serial. Prețul plătit
este mic: un salt suplimentar (datele mai trec o dată prin ESP32) și un pic de
suprasarcină de la protocolul de încadrare.

\begin{table}[H]
\centering
\caption{Bilanțul de cost orientativ al componentelor principale.}
\label{tab:bom}
\renewcommand{\arraystretch}{1.15}
\begin{tabular}{lc}
\toprule
\textbf{Componentă} & \textbf{Cost orientativ} \\
\midrule
Raspberry Pi 5 (4 GB)        & $\sim\SI{70}{\euro}$ \\
Senzor LIDAR LD14P           & $\sim\SI{50}{\euro}$ \\
ESP32 DevKit V1              & $\sim\SI{8}{\euro}$ \\
IMU MPU-6050                 & $\sim\SI{3}{\euro}$ \\
Tranzistor S8050, butoane, fire & $\sim\SI{5}{\euro}$ \\
Acumulator portabil USB (\textit{power bank}, 22,5 W) & $\sim\SI{30}{\euro}$ \\
\midrule
\textbf{Total estimat}      & $\sim\SI{166}{\euro}$ \\
\bottomrule
\end{tabular}
\end{table}

Cum se vede în Tabelul~\ref{tab:bom}, costul total se păstrează sub pragul de
\SI{200}{\euro} pe care mi-l propusesem. În secțiunile care urmează iau pe rând
fiecare componentă din figură.

\section{Raspberry Pi 5 — unitatea centrală de procesare}
\label{sec:rpi5}

Creierul aparatului este un \textbf{Raspberry~Pi~5} cu \SI{4}{} sau
\SI{8}{\giga\byte} RAM~\cite{rpi5brief}, cu procesor quad-core ARM Cortex-A76.
Rulează Ubuntu Server 24.04 LTS și duce, singur, întreaga stivă software: SLAM,
fuziunea senzorilor, serverul web, baza de date și punctul de acces WiFi. Nu
există un al doilea calculator — totul se întâmplă pe Pi.

Am ales Pi~5, și nu un model mai vechi, din două motive practice. În primul
rând, slam\_toolbox solicită serios procesorul la fiecare actualizare a hărții,
iar puterea suplimentară a Pi~5 ține latența sub control. În al doilea rând,
restul sistemului (containerul PostgreSQL, traficul USB) profită de
interfețele mai rapide ale plăcii. O parte din configurarea Pi-ului se face la
nivel de kernel, în scriptul de provizionare: se eliberează portul serial
hardware al plăcii (rămas, de altfel, neutilizat după ce LIDAR-ul a fost mutat
pe ESP32, dar păstrat ca variantă de rezervă).

Aparatul este alimentat de la un \textbf{acumulator portabil} (\textit{power
bank}) de \textbf{22,5~W}, conectat la mufa USB-C a Pi-ului. ESP32-ul — și,
prin el, LIDAR-ul — se alimentează la rândul lor din Pi, așa că totul merge de
la o singură sursă, fără niciun cablu la priză. Este exact ce am nevoie pentru
un aparat pe care îl porți în mână prin cameră.

\section{ESP32 DevKit V1 — microcontroler coprocesor}
\label{sec:esp32}

Lângă Pi stă un microcontroler \textbf{ESP32} (modul WROOM-32, pe o plăcuță
DevKit~V1). El are rolul de ,,om bun la toate" pentru partea de senzori: se
ocupă de IMU (prin I\textsuperscript{2}C), de cele trei butoane și de LIDAR
(îi citește ieșirea și îi comandă alimentarea). Practic, tot ce trebuie să
discute cu Pi-ul trece prin el.

Am ales ESP32 dintr-un cumul de motive simple. Plăcuța DevKit are deja o punte
USB-Serial la bord, deci legătura cu Pi-ul se face direct, fără hardware în
plus. Framework-ul Arduino~\cite{arduinoesp32} e matur și ușor de folosit, dar
îmi lasă totuși accesul la registre atunci când am nevoie. Cipul e mic și
foarte ieftin (în jur de \SI{8}{\euro}), are destui pini și, important pentru
mine, o a doua interfață serială hardware — de care a fost nevoie pentru
releul de date LIDAR. Firmware-ul nu folosește un sistem de operare în timp
real, ci un planificator simplu, cooperativ (§\ref{sec:scheduler}), care este
suficient pentru sarcina de aici.

\section{Senzorul LIDAR LD14P (LD-D200)}
\label{sec:lidar}

Senzorul propriu-zis este un \textbf{LDROBOT LD14P}~\cite{ld14p}: un mic
scaner laser care se rotește continuu și ,,mătură'' camera pe $360^\circ$.
Pentru fiecare unghi, măsoară distanța până la primul obstacol, cu ajutorul
unui fascicul laser în infraroșu și al unui fotodetector montate pe un ax
rotitor. Ce contează pentru proiect: are o rază utilă de până la
\SI{8}{\meter}, face în jur de zece rotații pe secundă, o rezoluție unghiulară
de aproape o jumătate de grad și trimite datele pe o linie serială la
\SI{230400}{} baud, cu nivel logic de \SI{3.3}{\volt} — deci se poate lega
direct la ESP32, fără adaptor de nivel.

Formatul în care senzorul își trimite datele nu este documentat oficial de
producător, însă există un driver open-source (\texttt{ldlidar\_stl\_ros2})
care îl înțelege și transformă fiecare rotație într-un mesaj standard de
scanare. Un detaliu care mi-a dat ceva bătăi de cap este că numărul de puncte
dintr-o rotație nu este fix, ci variază ușor (între 663 și 668) odată cu
turația motorului. Acest lucru a cerut o mică intervenție în driver, pe care o
descriu în §\ref{sec:ldlidar-driver}.

\subsection{Problema culorilor de fir non-standard}
\label{sec:lidar-wires}

Merită deschisă o paranteză despre cablajul LIDAR-ului, pentru că aici se
ascunde o capcană pe care am întâlnit-o direct. Cele patru fire ale senzorului
\textbf{nu} respectă convenția obișnuită de culori (la care negru = masă, roșu
= alimentare). La LD14P, conform fișei tehnice, lucrurile stau invers față de
ce te-ai aștepta: firul \textbf{negru} este alimentarea de \SI{5}{\volt},
\textbf{albul} duce datele, \textbf{roșul} este masa, iar \textbf{verdele} este
linia de recepție a senzorului (Tabelul~\ref{tab:lidar-wires}). În plus, firul
verde trebuie legat tot la masă, pentru ca senzorul să meargă la viteza lui
implicită — de aceea roșul și verdele ajung împreună în circuitul de comandă a
motorului (§\ref{sec:s8050}).

Insist asupra acestui detaliu fiindcă, dacă te iei după culorile standard,
fie nu pornești senzorul deloc, fie riști să îl arzi (legând din greșeală
alimentarea la masă). Eu însumi am inversat la început roșul cu verdele și
am pierdut timp până am verificat totul direct pe hardware.

\begin{table}[H]
\centering
\caption{Cablajul LD14P\\$\rightarrow$~ESP32 (note: culori non-standard).}
\label{tab:lidar-wires}
\begin{tabular}{lll}
\toprule
\textbf{Fir LD14P} & \textbf{Funcție} & \textbf{Conectează la} \\
\midrule
Negru  & VCC (5~V)  & ESP32 VIN \\
Alb    & TX (date)  & ESP32 GPIO~16 (Serial2 RX) \\
Verde  & RX         & Colector S8050 (cu Roșu) \\
Roșu   & GND        & Colector S8050 (cu Verde) \\
\bottomrule
\end{tabular}
\end{table}

\section{Unitatea inerțială MPU-6050}
\label{sec:mpu6050}

Ca să știe cum este orientat aparatul, folosesc un IMU \textbf{MPU-6050}%
~\cite{mpu6050} — un cip ieftin și foarte răspândit, care înglobează un
accelerometru și un giroscop, fiecare pe trei axe. Comunică cu ESP32-ul prin
magistrala I\textsuperscript{2}C (pinii GPIO~21 și 22), la \SI{400}{\kilo\hertz}.

L-am configurat (§\ref{sec:mpu-driver}) pe domenii potrivite pentru mișcarea
unei mâini: $\pm\SI{4}{\g}$ pe accelerometru și $\pm\SI{500}{\degree\per\second}$
pe giroscop, cu un filtru intern trece-jos la aproximativ \SI{44}{\hertz} și o
rată de ieșire de \SI{100}{\hertz}. Banda de \SI{44}{\hertz} acoperă lejer
mișcarea umană (care rar trece de vreo \SI{10}{\hertz}), dar taie din zgomotul
de vibrație de frecvență înaltă.

Exemplarul meu de MPU-6050 are însă o ciudățenie de care a trebuit să țin cont:
un bias de fabrică (\texttt{ZA\_OFFSET}) care face ca, în repaus, accelerometrul
să indice pe axa Z în jur de \SI{15.3}{\meter\per\second\squared} în loc de
\SI{9.81}{}~\cite{mpu6050regmap}. Cu alte cuvinte, sensul gravitației e corect,
dar mărimea e mult prea mare. Tentația firească ar fi să ,,corectez'' acest
offset scriind zero în registrele respective — dar am aflat, și din documentație,
și din încercări proprii, că nu trebuie făcut: unul dintre biți este rezervat
compensării de temperatură, iar o scriere naivă strică și mai rău calibrarea (am
ajuns la o gravitație de doar \SI{7.83}{\meter\per\second\squared}). Așa că las
cipul exact cum vine de la fabrică și las corecția pe seama algoritmilor de
fuziune. De aici vine, de altfel, și decizia de a ignora complet accelerația în
EKF (§\ref{sec:ekf}).

\section{Circuitul de control al motorului LIDAR (S8050)}
\label{sec:s8050}

Motorul LIDAR-ului nu trebuie să se rotească tot timpul — vreau să-l pot
porni și opri din software, ca să nu consume curent și să nu se uzeze degeaba
cât timp nu scanez. Pentru asta am folosit un tranzistor NPN ieftin
(\textbf{S8050}) pe post de întrerupător comandat, montat pe firul de masă al
senzorului (Figura~\ref{fig:s8050}). Pinul GPIO~4 al ESP32-ului îi comandă baza
printr-o rezistență de \SI{1}{\kilo\ohm}: când pinul e pe HIGH, tranzistorul
conduce și ,,închide'' masa LIDAR-ului, deci motorul pornește; când e pe LOW,
tranzistorul se blochează, LIDAR-ul rămâne fără masă și motorul se oprește.

Am preferat să comut masa, nu alimentarea, pentru că este și mai simplu, și mai
sigur. Mai simplu, fiindcă un NPN se comandă direct dintr-un pin de
\SI{3.3}{\volt}, fără componente suplimentare. Mai sigur, fiindcă, dacă
ESP32-ul rămâne nealimentat, tranzistorul se blochează de la sine, iar motorul
nu poate porni din greșeală.

Două detalii din firmware întăresc această siguranță. În primul rând, GPIO~4
este pus pe LOW chiar pe \textbf{prima linie} din \texttt{setup()}, înaintea
oricărei alte inițializări, ca motorul să stea oprit pe toată durata pornirii
ESP32-ului. În al doilea rând, există un \textbf{watchdog}: dacă Pi-ul nu mai
trimite, vreme de trei secunde, semnalul prin care cere motorul pornit,
firmware-ul îl oprește singur (§\ref{sec:watchdog}). Așa, dacă Pi-ul se
blochează sau cablul se desprinde, LIDAR-ul nu rămâne să se învârtă la
nesfârșit.

% Schema 2 — circuitul S8050
\begin{figure}[H]
\centering
\begin{tikzpicture}[
    every node/.style={font=\footnotesize},
    pin/.style={circle, draw, fill=black, inner sep=1pt}
]
    \node (gpio4) at (0, 2) {ESP32 GPIO~4};
    \node[pin] (a) at (1.5, 2) {};
    \draw (gpio4) -- (a);
    \node (res) at (3, 2) {1~k$\Omega$};
    \node[pin] (b) at (4.5, 2) {};
    \draw (a) -- (b);
    \node (base) at (5.5, 2) {Bază S8050};
    \draw (b) -- (base);

    \node (esp32gnd) at (0, 0.5) {ESP32 GND};
    \node[pin] (c) at (1.5, 0.5) {};
    \draw (esp32gnd) -- (c);
    \node (emitter) at (5.5, 0.5) {Emitor S8050};
    \draw (c) -- (emitter);

    \node (lidarred) at (0, -1) {LD14P Roșu};
    \node (lidargreen) at (0, -1.5) {LD14P Verde};
    \node[pin] (d) at (3, -1) {};
    \node[pin] (e) at (3, -1.5) {};
    \draw (lidarred) -- (d);
    \draw (lidargreen) -- (e);
    \draw (d) -- (e);
    \node (collector) at (5.5, -1.25) {Colector S8050};
    \draw (d) -- (5.5, -1) -- (collector);

    \node (vin) at (0, -3) {ESP32 5V (VIN)};
    \node (lidarblack) at (5.5, -3) {LD14P Negru};
    \draw (vin) -- (lidarblack);

    \node[right=0.5cm of base, font=\scriptsize, align=left] (note) {%
        $\mathrm{HIGH \Rightarrow}$ S8050 saturat $\Rightarrow$ motor pornit\\%
        $\mathrm{LOW}  \Rightarrow$ S8050 blocat $\Rightarrow$ motor oprit};
\end{tikzpicture}
\caption{Circuitul de control al motorului LIDAR cu tranzistor NPN S8050 ca \textit{low-side switch}. Pinul GPIO~4 al ESP32 controlează starea tranzistorului prin rezistența de bază de 1~k$\Omega$.}
\label{fig:s8050}
\end{figure}

\section{Butoanele și interfața utilizatorului}
\label{sec:buttons}

Aparatul are trei butoane fizice: \textbf{SHUTDOWN}, \textbf{RESET} și
\textbf{SAVE} (pe GPIO~25, 26 și 27). Sunt cele mai simple butoane cu putință:
fiecare leagă pinul la masă atunci când este apăsat, iar în rest pinul e ținut
sus de rezistența de pull-up internă a ESP32-ului. Practic, apăsarea înseamnă
,,pin tras la zero".

Nu am pus niciun circuit de dezamortizare (filtru, trigger Schmitt etc.) —
fac totul în software. Problema, cunoscută, este că un buton mecanic
,,zornăie" câteva milisecunde la apăsare, generând o serie de tranziții false.
Soluția din clasa \texttt{Button} (§\ref{sec:btn-driver}) este să aștepte ca
citirea să rămână stabilă \SI{25}{\milli\second} înainte de a o considera o
apăsare reală. Pe lângă apăsarea scurtă, codul recunoaște și apăsarea lungă
(peste \SI{2}{\second}), folosită pentru oprirea aparatului. Am ales această
abordare integral software pentru că nu costă nimic în componente, iar
întârzierea de \SI{25}{\milli\second} este complet imperceptibilă la
utilizare.

\clearpage

% =============================================================================
% CAPITOLUL 5 — PROTOCOLUL UART
% =============================================================================
\chapter{Protocolul de comunicație ESP32~\texorpdfstring{$\leftrightarrow$}{<->}~Raspberry Pi}
\label{cap:uart}

\section{Stratul fizic — UART prin USB-CDC}
\label{sec:uart-phy}

ESP32-ul și Raspberry~Pi vorbesc între ele printr-o legătură serială peste
USB: placa de dezvoltare are deja un cip care face conversia, așa că, din
punctul de vedere al Pi-ului, ESP32-ul apare ca un port serial obișnuit
(\texttt{/dev/ttyUSB0})~\cite{usbcdc}. Comunicația merge la \SI{460800}{} baud.

Inițial linia funcționa la \SI{115200}{} baud, ceea ce era de ajuns câtă vreme
prin ea treceau doar datele de la IMU. A trebuit însă să urc viteza odată ce am
început să trec și fluxul LIDAR prin ESP32. Socoteala e simplă: IMU-ul ocupă în
jur de \SI{3}{\kilo\byte\per\second}, iar fluxul LIDAR aproape
\SI{23}{\kilo\byte\per\second} — împreună, cam \SI{26}{\kilo\byte\per\second}.
La \SI{115200}{} baud (puțin peste \SI{11}{\kilo\byte\per\second}) doar fluxul
LIDAR ar fi depășit capacitatea legăturii, pe când \SI{460800}{} baud lasă o
rezervă confortabilă.

\section{Formatul cadrului binar}
\label{sec:frame-format}

% Schema 3 — formatul cadrului
\begin{figure}[H]
\centering
\begin{tikzpicture}[
    every node/.style={font=\footnotesize},
    box/.style={rectangle, draw, minimum height=0.8cm, align=center, font=\ttfamily}
]
    \node[box, minimum width=1.2cm] (s0) {0xA5};
    \node[box, minimum width=1.2cm, right=0pt of s0] (s1) {0x5A};
    \node[box, minimum width=1.5cm, right=0pt of s1, fill=blue!5] (type) {TYPE};
    \node[box, minimum width=1.5cm, right=0pt of type, fill=blue!5] (len) {LEN};
    \node[box, minimum width=2.5cm, right=0pt of len, fill=yellow!10] (pay) {PAYLOAD\\(0--64 B)};
    \node[box, minimum width=1.2cm, right=0pt of pay, fill=red!10] (crc) {CRC8};

    \node[below=0.1cm of s0, font=\tiny] {SYNC[0]};
    \node[below=0.1cm of s1, font=\tiny] {SYNC[1]};
    \node[below=0.1cm of type, font=\tiny] {1 B};
    \node[below=0.1cm of len, font=\tiny] {1 B};
    \node[below=0.1cm of pay, font=\tiny] {LEN B};
    \node[below=0.1cm of crc, font=\tiny] {1 B};

    \draw [decorate, decoration={brace, mirror, amplitude=5pt}]
        (type.south west) ++ (0,-0.6) -- ($(pay.south east) + (0,-0.6)$)
        node[midway, below=8pt] {acoperire CRC8};
\end{tikzpicture}
\caption{Format cadru binar UART. CRC-8 acoperă \textit{TYPE + LEN + PAYLOAD}; biții de sincronizare sunt excluși pentru a permite resincronizare independentă.}
\label{fig:frame-format}
\end{figure}

Datele nu circulă pur și simplu ca un șir de octeți, ci grupate în
\textit{cadre}, ca să se știe unde începe și unde se termină fiecare mesaj.
Structura unui cadru este cea din Figura~\ref{fig:frame-format}: doi octeți de
sincronizare ($\mathtt{0xA5}$ și $\mathtt{0x5A}$) care marchează începutul,
apoi un octet pentru tipul mesajului, unul pentru lungimea conținutului,
conținutul propriu-zis și, la final, un octet de control (CRC). Numerele de
mai mulți octeți sunt scrise în ordine \textit{little-endian} pe ambele
capete, așa că pot fi copiate direct, fără conversii.

Suma de control acoperă tipul, lungimea și conținutul, dar \textbf{nu} și
octeții de sincronizare. I-am lăsat dinadins pe dinafară, ca să pot căuta
începutul unui cadru (§\ref{sec:resync}) independent de verificarea CRC.

Cei doi octeți de sincronizare nu sunt aleși la întâmplare. Combinația
$\mathtt{0xA5\,0x5A}$ este greu de întâlnit din întâmplare într-un flux de
valori de la IMU și, fiind antisimetrică pe biți, reduce și șansa unei false
sincronizări dacă fluxul ar fi citit decalat cu un bit.

\section{Tipuri de cadre}
\label{sec:frame-types}

\begin{table}[H]
\centering
\caption{Tipurile de cadru implementate în firmware (\texttt{firmware/esp32/src/config.h}).}
\label{tab:frame-types}
\renewcommand{\arraystretch}{1.15}
\begin{tabular}{llcll}
\toprule
\textbf{Hex} & \textbf{Nume} & \textbf{Direcție} & \textbf{LEN} & \textbf{Conținut} \\
\midrule
0x01 & IMU         & ESP$\to$Pi & 24    & 6 $\times$ float32 (m/s², rad/s) \\
0x02 & BUTTON      & ESP$\to$Pi & 2     & uint8 id, uint8 state \\
0x03 & HEARTBEAT   & ESP$\to$Pi & 4     & uint32 uptime\_ms \\
0x04 & STATUS      & ESP$\to$Pi & 2 / 8 & flags + diagnostic chip \\
0x05 & LIDAR\_FRAME & ESP$\to$Pi & 1--64 & octeți bruti LD14P \\
0x06 & LIDAR\_EN   & Pi$\to$ESP & 1     & uint8 (0=off, 1=on) \\
0x07 & LIDAR\_ACK  & ESP$\to$Pi & 1     & uint8 (motor state) \\
\bottomrule
\end{tabular}
\end{table}

\section{Calculul CRC-8}
\label{sec:crc-calc}

Suma de control se calculează la fel pe ambele capete — o dată în C++, pe
ESP32, și o dată în Python, pe Pi. Algoritmul propriu-zis este cel clasic, pe
biți: pentru fiecare octet primit se actualizează o valoare de 8 biți printr-o
succesiune de operații de tip XOR și deplasare, folosind polinomul
$\mathtt{0x07}$.

Partea importantă nu este atât formula, cât faptul că am ales exact aceeași
variantă de CRC ca utilitarul \texttt{crc8} din nucleul Linux~\cite{linuxcrc8,
maxim1wirecrc}. Așa am putut verifica ușor că cele două implementări (cea de pe
ESP32 și cea de pe Pi) dau rezultate identice — câteva valori de probă folosite
la testare sunt în Tabelul~\ref{tab:crc-vectors}.

\begin{table}[H]
\centering
\caption{Vectori de test pentru CRC-8 (poli.\ $\mathtt{0x07}$, init $\mathtt{0x00}$).}
\label{tab:crc-vectors}
\renewcommand{\arraystretch}{1.15}
\begin{tabular}{ll}
\toprule
\textbf{Intrare} & \textbf{CRC-8} \\
\midrule
\texttt{b""}                  & $\mathtt{0x00}$ \\
\texttt{b"\textbackslash x01"} & $\mathtt{0x07}$ \\
\texttt{b"hello"}             & $\mathtt{0x92}$ \\
\bottomrule
\end{tabular}
\end{table}

\section{Algoritmul de resincronizare a parserului}
\label{sec:resync}

Pe o legătură serială se mai pierd sau se strică octeți, așa că cel care
citește cadrele (și pe Pi, și în firmware) trebuie să-și poată regăsi singur
ritmul. L-am scris ca o mică mașină de stări care, în esență, parcurge mereu
aceiași pași: întâi caută cei doi octeți de sincronizare, apoi citește tipul,
lungimea și conținutul, iar la final verifică suma de control.

Două amănunte o fac rezistentă la erori. Dacă, în timp ce caută al doilea
octet de sincronizare, dă peste încă un $\mathtt{0xA5}$, nu se zăpăcește, ci
rămâne în așteptare — astfel tratează corect cazul în care primul $\mathtt{0xA5}$
era doar zgomot, iar cadrul adevărat începe de la al doilea. Iar dacă lungimea
anunțată în cadru depășește maximul permis (semn că un octet rătăcit a fost luat
drept câmp de lungime), aruncă pur și simplu cadrul și o ia de la capăt. Toate
aceste situații — octet de sincronizare pierdut, sumă de control greșită,
lungime aberantă sau date sosite fragmentat — sunt verificate de testele
automate ale parserului.

\section{Watchdog-ul motorului LIDAR}
\label{sec:watchdog}

Comanda motorului vine de la Pi: cât timp vrea motorul pornit, Pi-ul trimite
spre ESP32, o dată pe secundă, un semnal de tip ,,ține motorul aprins". La
fiecare astfel de semnal, firmware-ul reîncarcă un cronometru de siguranță
(watchdog). Dacă trec trei secunde fără niciun semnal, firmware-ul presupune că
ceva nu e în regulă și oprește motorul singur. Pragul de trei secunde este
ales suficient de mare peste intervalul de o secundă, ca o mică întârziere de
rețea să nu ducă la opriri nedorite.

Toate comparațiile de timp bazate pe \texttt{millis()} din firmware folosesc
aritmetică cu semn. Aceasta evită o eroare clasică pe microcontrolere: o
scădere fără semn între două citiri de timp foarte apropiate poate produce,
prin depășire, un număr foarte mare care ar declanșa watchdog-ul în mod fals.
Cu aritmetică cu semn diferența rămâne corectă, iar motorul se oprește doar
atunci când chiar lipsesc reîmprospătările de la Raspberry Pi.

\clearpage

% =============================================================================
% CAPITOLUL 6 — FIRMWARE ESP32
% =============================================================================
\chapter{Firmware-ul ESP32}
\label{cap:firmware}

\section{Mediul de dezvoltare PlatformIO}
\label{sec:pio}

Firmware-ul de pe ESP32 l-am scris și compilat cu \textbf{PlatformIO}%
~\cite{platformio}, peste framework-ul Arduino~\cite{arduinoesp32}. Am
preferat PlatformIO în locul mediului Arduino clasic pentru că lucrează din
linia de comandă și gestionează singur dependențele și \textit{toolchain}-ul,
ceea ce face compilarea ușor de repetat. Codul este în C++.

Un lucru pe care l-am urmărit a fost să nu depind de biblioteci externe: în
afară de \texttt{Arduino.h} și de cea pentru magistrala I\textsuperscript{2}C,
nu folosesc nimic altceva. Aș fi putut lua o bibliotecă gata făcută pentru
MPU-6050, dar am preferat să scriu eu citirile, ca să am control complet
asupra inițializării senzorului. Per total, firmware-ul ocupă în jur de
22\% din memoria flash a cipului, deci rămâne loc berechet pentru eventuale
adăugiri.

\section{Planificatorul cooperativ}
\label{sec:scheduler}

Firmware-ul nu folosește un sistem de operare în timp real (de tip FreeRTOS),
ci o buclă simplă, \textbf{cooperativă}. Funcția \texttt{loop()} se execută cât
de des poate, iar fiecare sarcină (citirea IMU, trimiterea heartbeat-ului,
verificarea butoanelor etc.) se execută doar când îi vine ,,rândul", adică
atunci când a trecut destul timp de la ultima dată. Ordinea pașilor dintr-o
iterație este cea din Figura~\ref{fig:loop-firmware}: se citesc cadrele venite
de la Pi, se transmit datele LIDAR (dacă motorul e pornit), se verifică
watchdog-ul și butoanele, apoi se eșantionează IMU-ul și se trimite
heartbeat-ul.

Am ales abordarea simplă pentru că sarcina nu cere ceva mai complicat: bucla se
execută de peste o mie de ori pe secundă, mult mai des decât au nevoie cei
\SI{100}{\hertz} ai IMU-ului. Un sistem cu task-uri preemptive ar fi adus în
plus mai mult complexitate și riscuri (sincronizare între task-uri, inversiune
de prioritate) decât beneficii reale aici~\cite{li2003realtime, barr1999embedded}.
Singura grijă este ca niciun pas să nu blocheze bucla; de aceea, de pildă,
citirile de pe magistrala IMU au un timp maxim de așteptare, ca o defecțiune a
senzorului să nu înțepenească totul.

% Schema 4 — fluxul loop() al firmware-ului
\begin{figure}[H]
\centering
\begin{tikzpicture}[
    node distance=0.4cm,
    every node/.style={font=\scriptsize},
    block/.style={rectangle, draw, rounded corners, minimum width=4.5cm, minimum height=0.6cm, align=center},
    arr/.style={-{Latex[length=2mm]}}
]
    \node[block, fill=blue!5]   (n1) {now = millis()};
    \node[block, below=of n1]   (n2) {drain Serial (cadre Pi$\to$ESP)};
    \node[block, below=of n2]   (n3) {if g\_lidar\_enabled: drain Serial2 (LIDAR)};
    \node[block, below=of n3]   (n4) {watchdog LIDAR (signed compare)};
    \node[block, below=of n4]   (n5) {buton SHUTDOWN/RESET/SAVE.update()};
    \node[block, below=of n5]   (n6) {if now $\geq$ next\_imu\_ms: imu::read() + send\_imu};
    \node[block, below=of n6]   (n7) {if now $\geq$ next\_heartbeat\_ms: send\_heartbeat};
    \node[block, below=of n7]   (n8) {update\_status\_led()};

    \draw[arr] (n1) -- (n2);
    \draw[arr] (n2) -- (n3);
    \draw[arr] (n3) -- (n4);
    \draw[arr] (n4) -- (n5);
    \draw[arr] (n5) -- (n6);
    \draw[arr] (n6) -- (n7);
    \draw[arr] (n7) -- (n8);
    \draw[arr, dashed] (n8.east) .. controls +(1.2,0) and +(1.2,0) .. (n1.east);
\end{tikzpicture}
\caption{Bucla cooperativă \texttt{loop()} a firmware-ului ESP32. Fiecare iterație drenează cozile de I/O și verifică deadline-urile bazate pe \texttt{millis()}.}
\label{fig:loop-firmware}
\end{figure}

\section{Driverul MPU-6050}
\label{sec:mpu-driver}

Driverul pentru MPU-6050 (\texttt{imu.cpp}) comunică direct cu cipul, la nivel
de registru, și este scris cu grijă față de câteva capcane cunoscute ale acestui
senzor.

La pornire, driverul verifică mai întâi că răspunde chiar un MPU-6050, apoi îi
face un \textbf{reset complet} și așteaptă o zecime de secundă. Pasul acesta de
așteptare nu este de prisos: dacă scrii setările prea devreme după alimentare,
cipul le ignoră pe tăcute, iar accelerometrul rămâne pe domeniul implicit —
simptomul fiind că gravitația apare de două ori mai mare decât ar trebui. Abia
după reset îi configurez domeniile alese ($\pm\SI{4}{\g}$ și
$\pm\SI{500}{\degree\per\second}$), filtrul intern și rata de \SI{100}{\hertz}.
La final, recitesc registrele de configurare ca să mă asigur că setările chiar
au ,,prins"; dacă nu, raportez eroare în loc să trimit mai departe date greșite.

Transformarea valorilor brute ale cipului în unități fizice o fac tot pe ESP32,
cu doi factori de scalare ficși:
\[
\texttt{ACCEL\_LSB\_TO\_MS2} = \frac{9{,}80665}{8192}, \qquad
\texttt{GYRO\_LSB\_TO\_RADS} = \frac{\pi/180}{65{,}5}.
\]
O citire aduce, într-o singură tranzacție, toate cele șase valori (plus
temperatura, pe care o ignor). Aici intervine și capcana \texttt{ZA\_OFFSET}
de care vorbeam în §\ref{sec:mpu6050}: registrele de offset sunt doar citite,
niciodată scrise, fiindcă o scriere naivă ar strica și mai rău calibrarea.

\section{Driverul butoanelor cu debounce}
\label{sec:btn-driver}

Dezamortizarea butoanelor, despre care am vorbit ca principiu în
§\ref{sec:buttons}, este implementată într-o clasă \texttt{Button}, câte una
pentru fiecare buton. Logica este cea pe care te-ai aștepta: clasa ține minte
starea ,,stabilă" curentă și momentul ultimei schimbări, iar o apăsare sau o
eliberare este luată în seamă doar după ce semnalul rămâne neschimbat
\SI{25}{\milli\second}. Pe lângă apăsarea scurtă, clasa recunoaște și apăsarea
lungă (peste \SI{2}{\second}), folosită pentru oprirea aparatului. De fiecare
dată când se petrece ceva — apăsare, eliberare sau apăsare lungă — clasa
anunță restul firmware-ului printr-o funcție de \textit{callback}, iar de acolo
evenimentul pleacă mai departe spre Pi.

\section{Releul datelor LIDAR}
\label{sec:lidar-relay}

Releul de date LIDAR este partea din firmware care face din ESP32 un simplu
repetor: tot ce primește de la senzor pe a doua interfață serială împachetează
și trimite mai departe către Pi, fără să se uite la conținut. Octeții ajung pe
Pi în aceeași ordine, iar acolo sunt redați driverului de LIDAR.

Singura decizie care a cerut puțină atenție a fost dimensionarea memoriei
tampon de recepție. Fluxul LIDAR este destul de rapid
(\textasciitilde\SI{23}{\kilo\byte\per\second}), iar tamponul implicit s-ar
umple în vreo \SI{10}{\milli\second} — un interval pe care o iterație de buclă
mai lentă l-ar putea depăși din când în când, pierzând astfel o bucată de
scanare. De aceea am mărit tamponul la \SI{1}{\kilo\byte}, care se umple abia
în vreo \SI{45}{\milli\second}, mult peste durata unei iterații normale.
Releul lucrează numai cât timp motorul este pornit; cu motorul oprit, orice
octet rămas pe linie este doar zgomot și este ignorat.

\clearpage

% =============================================================================
% CAPITOLUL 7 — STIVA ROS2 PE PI
% =============================================================================
\chapter{Stiva software pe Raspberry Pi (ROS2)}
\label{cap:ros2}

\section{ROS2 Jazzy Jalisco — cadrul de lucru}
\label{sec:ros2-overview}

Toată partea de calcul de pe Pi este construită peste \textbf{ROS2 Jazzy
Jalisco}~\cite{ros2jazzy}. ROS2 este, pe scurt, un cadru de lucru pentru
robotică: îți oferă o mulțime de ,,piese" gata făcute și un mod standard prin
care ele schimbă date între ele. Am ales ROS2 (și nu ROS1, varianta mai veche)
pentru că este versiunea actuală, întreținută activ, iar pachetele de care
aveam nevoie — slam\_toolbox și robot\_localization — există direct pentru
ea~\cite{maruyama2016ros2, ros2why, ros2macenski}.

Câteva noțiuni din ROS2 revin des în continuare, așa că le explic pe scurt.
Un sistem ROS2 este alcătuit din \textbf{noduri} (programe care rulează separat),
care comunică între ele prin \textbf{topice} (un fel de canale pe care un nod
publică, iar altele citesc — de exemplu \texttt{/scan} sau \texttt{/map}) și
prin \textbf{servicii} (apeluri de tip cerere-răspuns). Pe lângă acestea,
există \textbf{arborele de transformări} (TF2), care ține evidența relațiilor
geometrice dintre diferitele sisteme de coordonate ale sistemului. Cum sunt
legate acestea în proiect se vede în Figura~\ref{fig:tf-tree}. Un avantaj al
acestei organizări este că pot înlocui un senzor real cu unul simulat fără să
schimb nimic în restul sistemului.

% Schema 5 — arborele TF
\begin{figure}[H]
\centering
\begin{tikzpicture}[
    every node/.style={font=\footnotesize},
    frame/.style={rectangle, draw, rounded corners, minimum width=1.8cm, minimum height=0.7cm, align=center, fill=blue!5},
    arr/.style={-{Latex[length=2mm]}, thick}
]
    \node[frame] (map) {map};
    \node[frame, below=1cm of map] (odom) {odom};
    \node[frame, below=1cm of odom] (base) {base\_link};
    \node[frame, below right=0.6cm and 0.3cm of base] (laser) {laser\_frame};
    \node[frame, below left=0.6cm and 0.3cm of base] (imu) {imu\_link};

    \draw[arr] (map) -- node[right]{\scriptsize \texttt{slam\_toolbox} (dinamic, $\sim$50 Hz)} (odom);
    \draw[arr] (odom) -- node[right]{\scriptsize \texttt{ekf\_node} (dinamic, 30 Hz)} (base);
    \draw[arr] (base) -- node[right]{\scriptsize static (0, 0, 0.10 m)} (laser);
    \draw[arr] (base) -- node[left]{\scriptsize static (identitate)} (imu);
\end{tikzpicture}
\caption{Arborele de transformări TF2 din timpul rulării în modul \texttt{setup.sh full}. \texttt{slam\_toolbox} publică \texttt{map}~$\to$~\texttt{odom}; \texttt{robot\_localization} publică \texttt{odom}~$\to$~\texttt{base\_link}.}
\label{fig:tf-tree}
\end{figure}

\section{Pachetul recon\_hardware}
\label{sec:recon-hardware}

Pachetul \texttt{recon\_hardware} cuprinde nodurile care leagă restul
sistemului de lumea fizică. Cel mai important este puntea către ESP32
(\texttt{esp32\_uart\_bridge.py}). Ea citește fluxul de cadre binare de pe
portul serial, le desface și trimite fiecare tip de informație pe topicul ROS
potrivit: datele IMU pe \texttt{/imu/data\_raw}, apăsările de buton pe
topicele \texttt{/buttons/\ldots}, iar o dată pe secundă un mic raport de stare
a legăturii. Tot această punte oferă și serviciul prin care Pi-ul pornește și
oprește motorul LIDAR, repetând comanda de pornire o dată pe secundă (semnalul
pe care îl așteaptă watchdog-ul din firmware).

Citirea portului serial se face pe un fir de execuție separat, ca apelurile
blocante să nu încetinească restul nodului. Datele LIDAR primite sunt scrise
de punte într-un \textbf{pseudo-terminal}~\cite{linuxpty} --- un port serial
„virtual" oferit de Linux. Driverul LD14P deschide acest pseudo-terminal exact
ca pe un port serial real și nu „observă" că octeții vin, de fapt, prin ESP32.
Acesta este artificiul care permite reutilizarea driverului de LIDAR nemodificat,
fără o conexiune directă la UART-ul Raspberry Pi.

Al doilea nod din pachet preia datele brute de la giroscop și calculează din
ele unghiul de orientare, după metoda descrisă în §\ref{sec:giro}. L-am
construit în jurul unei funcții ,,pure" (una care depinde doar de intrările ei,
fără stare ascunsă), tocmai ca să o pot testa ușor și separat de restul
sistemului.

\section{Pachetul recon\_db}
\label{sec:recon-db}

Pachetul \texttt{recon\_db} se ocupă de salvarea hărților. Nodul său ascultă
în permanență harta curentă și o ține pregătită în memorie, ca o salvare să se
poată face instantaneu. Când primește comanda de salvare — fie de la butonul
fizic SAVE de pe aparat, fie din interfața web — scrie harta în baza de date,
cu un nume bazat pe dată și oră, și notează evenimentul într-un jurnal pe care
interfața web îl verifică periodic.

Operațiile cu baza de date se execută pe fire de execuție separate, ca să nu
blocheze nodul cât timp scrie pe disc. Salvarea poate fi pornită și fără ca
interfața web să fie deschisă, doar de la butonul fizic — o facilitate
moștenită din varianta inițială a proiectului, pe care am păstrat-o.

\section{slam\_toolbox — cartografierea}
\label{sec:slam-tb}

Cartografierea propriu-zisă o face slam\_toolbox, rulat în mod \textit{online}
(asincron): procesează fiecare scanare imediat ce sosește, fără să aștepte,
exact ce trebuie pentru funcționare în timp real pe Pi. Câteva dintre setările
din fișierul de configurare merită explicate:

\begin{itemize}\itemsep0pt
    \item harta este reconstruită \textbf{în fiecare secundă}; un interval
          scurt face ca obstacolele trecătoare (de pildă o persoană care
          tocmai a trecut prin cameră) să dispară repede din hartă, cu prețul
          unui consum ceva mai mare de procesor;
    \item potrivirea de scanări este lăsată să ruleze des, ca să țină pasul cu
          mișcarea mâinii;
    \item am cerut explicit ca slam\_toolbox să publice poziția \textbf{chiar
          și când aparatul stă pe loc} — altfel interfața web nu ar mai ști
          unde se află dispozitivul;
    \item în fine, i-am dat ca reper orientarea venită de la giroscop, ca
          ajutor în potrivirea scanărilor.
\end{itemize}

Tot la integrarea cu slam\_toolbox am dat peste cea mai supărătoare problemă
din proiect. Karto reține numărul de puncte din prima scanare pe care o
primește și \textbf{respinge} apoi orice scanare cu un număr diferit. Iar
LD14P-ul, după cum am spus, trimite de fiecare dată un număr ușor diferit de
puncte. Rezultatul era că harta se construia pentru o singură rotație și apoi
se oprea — fără niciun mesaj de eroare evident. Mi-a luat ceva timp până am
prins cauza; soluția a fost o mică modificare în driverul de LIDAR, care aduce
fiecare rotație la același număr fix de puncte (o descriu în
§\ref{sec:ldlidar-driver}).

\section{robot\_localization — filtrul Kalman extins}
\label{sec:ekf-node}

Filtrul Kalman extins, despre care am discutat teoretic în §\ref{sec:ekf},
este pus în practică tot printr-un pachet gata făcut
(robot\_localization~\cite{moore2014robotlocalization}). L-am configurat în mod
2D, la \SI{30}{\hertz}, și l-am pus să publice transformarea
\texttt{odom}~$\to$~\texttt{base\_link} (cea care, în versiunile anterioare,
era fixată artificial la identitate). Ca intrare primește orientarea de la
IMU.

Filtrul ar putea, în principiu, să fuzioneze toate cele 15 mărimi posibile —
poziție, orientare, viteze și accelerații pe cele trei axe:
\[
[\,x,\; y,\; z,\; \mathrm{roll},\; \mathrm{pitch},\; \mathrm{yaw},\;
\dot{x},\; \dot{y},\; \dot{z},\; \dot{\mathrm{roll}},\; \dot{\mathrm{pitch}},\;
\dot{\mathrm{yaw}},\; \ddot{x},\; \ddot{y},\; \ddot{z}\,].
\]
Eu însă i-am activat \textbf{doar} unghiul de cap și viteza lui de variație.
Restul sunt oprite, iar accelerația — în mod special — din motivul deja
cunoscut: biasul \texttt{ZA\_OFFSET} al cipului (§\ref{sec:mpu6050}) ar fi
integrat de două ori și ar duce la o derivă necontrolată. Poziția aparatului
rămâne, ca și până acum, în grija corecției venite de la SLAM.

\section{Driverul vendorat ldlidar\_stl\_ros2}
\label{sec:ldlidar-driver}

Pentru a citi senzorul am pornit de la driverul open-source al producătorului
(\texttt{ldlidar\_stl\_ros2}), pe care a trebuit însă să îl modific în două
locuri. Prima modificare îl face să citească de pe pseudo-terminalul creat de
punte (nu de pe portul serial al Pi-ului) și îl lasă să aștepte oricât primul
pachet de date — altfel, driverul renunța dacă nu primea nimic în trei secunde,
exact situația în care motorul este încă oprit. A doua modificare este cea care
rezolvă blocarea hărții descrisă mai sus: aduce fiecare rotație la un număr fix
de 720 de puncte.

Aceste modificări vin însă cu o capcană care merită semnalată: driverul este
inclus în proiect ca un sub-depozit git, dar fără configurarea care l-ar aduce
automat la clonare. Practic, dacă cineva clonează proiectul ,,de-a gata",
folderul driverului rămâne gol, modificările mele se pierd, iar harta se va
opri din nou după o rotație. Pe viitor, driverul ar trebui fie inclus ca
submodul propriu-zis, fie copiat direct în arborele proiectului.

\clearpage

% =============================================================================
% CAPITOLUL 8 — INTERFAȚA WEB
% =============================================================================
\chapter{Interfața web (recon\_webui)}
\label{cap:webui}

\section{Arhitectura Flask + SocketIO + eventlet}
\label{sec:flask-arch}

Interfața cu utilizatorul este o pagină web servită chiar de pe Raspberry~Pi.
Am construit-o cu \textbf{Flask}~\cite{flask}, la care am adăugat
\textbf{SocketIO}~\cite{flasksocketio, socketio} pentru a putea trimite date
către browser în timp real (de exemplu harta, pe măsură ce se construiește),
nu doar la cerere. În spate, serverul folosește biblioteca
\textbf{eventlet}~\cite{eventlet}, care îi permite să țină multe conexiuni
deschise în paralel fără să consume multe resurse — important pe un calculator
modest cum este Pi-ul.

Acest mod de lucru vine cu o regulă strictă: o anumită linie de inițializare a
bibliotecii eventlet trebuie să fie chiar prima executată în program, înaintea
oricărui alt import. Motivul este că eventlet ,,rescrie" la pornire o parte din
funcțiile standard ale limbajului ca să le facă cooperative; dacă unele dintre
ele sunt deja folosite înainte de această rescriere, apar blocaje greu de
depistat. De aceea ordinea importurilor din fișierul principal este fixă și
marcată ca atare în cod.

\section{Modelul de threading și sincronizare}
\label{sec:threading}

% Schema 6 — modelul de threading
\begin{figure}[H]
\centering
\begin{tikzpicture}[
    every node/.style={font=\scriptsize},
    hub/.style={rectangle, draw, rounded corners, fill=green!10, minimum width=6cm, minimum height=2cm, align=center},
    th/.style={rectangle, draw, rounded corners, fill=orange!10, minimum width=6cm, minimum height=1.5cm, align=center},
    arr/.style={-{Latex[length=2mm]}, thick}
]
    \node[hub] (hub) {%
        \textbf{Hub eventlet (greenlets)}\\
        \begin{tabular}{l}
        $\bullet$ rute Flask (toate cererile HTTP)\\
        $\bullet$ emit\_loop (SocketIO push)\\
        $\bullet$ interogări DB prin \texttt{tpool.execute(...)}\\
        $\bullet$ apeluri service ROS2 (poll \texttt{Future} cu \texttt{time.sleep})\\
        \end{tabular}};
    \node[th, below=0.6cm of hub] (ros) {%
        \textbf{Thread OS real ,,ros2\_bridge\_spin"}\\
        \texttt{rclpy.spin(node)} — callback-uri subscriber\\
        Publică în \texttt{queue.Queue(maxsize=64)}};
    \node[th, below=0.4cm of ros] (ser) {%
        \textbf{Thread OS real ,,esp32\_serial\_reader"}\\
        Citește \texttt{/dev/ttyUSB0}, parsează frame-uri\\
        Publică direct pe topice ROS};

    \draw[arr] (ros) -- node[right, align=center]{\texttt{queue.Queue}\\(evenimente)} (hub);
\end{tikzpicture}
\caption{Modelul de threading al procesului \texttt{recon\_webui}. Trei contexte de execuție coexistă: hub-ul eventlet (verde), thread-ul de \texttt{rclpy.spin}, și thread-ul de citire serial.}
\label{fig:threading}
\end{figure}

Procesul \texttt{recon\_webui} rulează în paralel trei contexte de execuție
(Figura~\ref{fig:threading}): bucla cooperativă eventlet (unde trăiesc rutele
Flask și emiterea pe WebSocket), un fir dedicat pentru \texttt{rclpy.spin()}
(ascultarea topicelor ROS2) și un fir pentru citirea portului serial.

Regula care menține sistemul stabil este simplă: mesajele către clienții web
(\texttt{socketio.emit()}) se trimit numai din bucla eventlet, niciodată din
firul ROS2. Datele primite de pe topice sunt trecute spre buclă printr-o coadă
(\texttt{queue.Queue}), iar interogările de bază de date din rutele web se
execută într-un fir separat (\texttt{tpool.execute()}) ca să nu blocheze bucla.
Combinația eventlet + ROS2 în același proces a generat, în timpul dezvoltării,
blocaje greu de depistat; odată respectată această separare a firelor,
comportamentul a devenit stabil.

\section{Pattern-ul DataChannel cu fallback automat}
\label{sec:datachannel}

Un detaliu de proiectare de care sunt mulțumit este că interfața nu are cod
separat pentru ,,demonstrație" și pentru ,,funcționare reală". Fiecare sursă de
date (poza, harta, valorile de la IMU și așa mai departe) trece printr-un mic
obiect care decide singur ce să afișeze: dacă a primit recent date reale de la
senzori, le folosește pe acelea; dacă nu, întoarce automat niște date simulate.
Așa, aceeași interfață merge și fără hardware (util pentru dezvoltare și
testare), și cu senzorii reali, fără nicio modificare de cod. Singura diferență
vizibilă este un filigran ,,MOCK", afișat atunci când datele sunt simulate.

Cum \texttt{DataChannel} este citit atât din bucla eventlet, cât și din firul
ROS2, accesul la valoarea curentă este protejat printr-un mecanism de blocare
(\textit{lock}), astfel încât cele două fire să nu intre în conflict când unul
scrie și celălalt citește.

\section{Paginile interfeței}
\label{sec:pages}

Interfața are trei pagini. Prima și cea mai folosită este \textbf{harta live}:
pe ea se desenează harta pe măsură ce se construiește, împreună cu poziția
curentă a aparatului și cu traseul parcurs (ultimele câteva sute de poziții,
sub forma unei dâre albastre). Tot de aici se pornește și se oprește scanarea
și se poate salva sau șterge harta, inclusiv de la tastatură. A doua pagină
este o \textbf{galerie} a hărților salvate, unde fiecare hartă poate fi privită
fie în forma brută, fie procesată (cu pereții scoși în evidență). A treia,
pagina de \textbf{diagnoză}, arată în timp real valorile de la IMU, starea
legăturii cu ESP32-ul și câteva statistici despre hartă.

Harta este desenată mereu \textbf{centrată pe poziția curentă}: aparatul
rămâne în mijlocul ecranului, iar harta ,,alunecă" în jurul lui pe măsură ce
te miști prin cameră. Ca pagina să funcționeze și fără acces la internet
(dispozitivul fiind pe rețeaua lui izolată), toate resursele sunt incluse
local, fără a depinde de servere externe.

\section{Punctul de acces WiFi}
\label{sec:wifi-ap}

Ca dispozitivul să fie de sine stătător, fără să depindă de o rețea
existentă, Raspberry~Pi își creează propria rețea WiFi. Folosesc pentru asta o
interfață virtuală dedicată punctului de acces, astfel încât placa de rețea
reală să rămână liberă, în caz că vreau totuși să conectez Pi-ul, în paralel,
la o rețea obișnuită.

De partea practică se ocupă două programe consacrate: unul
(\textbf{hostapd}~\cite{hostapd}) creează rețeaua WiFi protejată cu parolă
(WPA2), iar celălalt (\textbf{dnsmasq}~\cite{dnsmasq}) atribuie automat adrese
dispozitivelor care se conectează și permite accesul la interfață printr-o
adresă ușor de reținut, de tipul \texttt{recon.local}. Toată această configurare
se face automat la instalare și pornește singură odată cu aparatul. Trebuie
spus că interfața web nu cere parolă — ceea ce este acceptabil cât timp totul
rămâne pe această rețea izolată, dar care ar trebui adăugat înainte de a expune
dispozitivul pe o rețea publică (vezi §\ref{sec:limitations}).

\clearpage

% =============================================================================
% CAPITOLUL 9 — POSTGRESQL + TIER 2
% =============================================================================
\chapter{Stratul de persistență și post-procesarea hărților}
\label{cap:db}

\section{Containerul Docker PostgreSQL}
\label{sec:postgres-docker}

Hărțile salvate sunt păstrate într-o bază de date \textbf{PostgreSQL}%
~\cite{postgresql}, pe care o rulez într-un container \textbf{Docker}%
~\cite{docker}. Am ales această variantă pentru că îmi izolează baza de date
de restul sistemului și o face ușor de pornit și de mutat. Datele propriu-zise
stau într-un volum separat, deci rămân pe loc chiar dacă recreez containerul,
iar \texttt{setup.sh} așteaptă ca baza să fie gata înainte de a porni nodurile
care depind de ea. Baza de date este accesibilă doar local, de pe Pi.

Un amănunt: containerul și utilizatorul bazei poartă încă numele
\texttt{roomba}, moștenit din varianta inițială a proiectului. Le-am lăsat
intenționat așa, pentru a nu pierde scanările deja salvate; aplicația oricum nu
ține cont de aceste nume, ci se conectează prin intermediul unei variabile de
configurare.

\section{Schema bazei de date}
\label{sec:db-schema}

Baza de date are patru tabele, rezumate în Tabelul~\ref{tab:db-schema}. Le-am
descris în cod cu ajutorul unei biblioteci de tip ORM~\cite{sqlalchemy}, care
îmi permite să lucrez cu tabelele ca și cum ar fi obiecte, fără să scriu SQL de
mână. Pe scurt, există un tabel pentru hărțile salvate, unul pentru sesiunile
de lucru, unul pentru hărțile procesate (legat de harta-sursă, astfel încât
ștergerea unei hărți elimină automat și procesările ei) și unul care ține un
jurnal al salvărilor și ștergerilor, consultat de interfața web.

\begin{table}[H]
\centering
\caption{Tabelele schemei de bază de date.}
\label{tab:db-schema}
\renewcommand{\arraystretch}{1.15}
\begin{tabular}{ll}
\toprule
\textbf{Tabel} & \textbf{Rol} \\
\midrule
\texttt{maps}           & grila de ocupare salvată + metadate geometrie \\
\texttt{sessions}       & sesiuni de operare (FK opțional $\to$ \texttt{maps}) \\
\texttt{processed\_maps}& rezultate Tier-2 (FK CASCADE $\to$ \texttt{maps}) \\
\texttt{map\_events}    & jurnal SAVED/DELETED pentru \textit{polling} UI \\
\bottomrule
\end{tabular}
\end{table}

Harta propriu-zisă o salvez în format JSON, nu într-un format binar compact.
Am preferat JSON-ul pentru că este lizibil — pot inspecta o hartă direct în
baza de date — și, mai ales, pentru că îmi permite să adaug câmpuri noi (de
pildă segmentele de perete sau etichetele clusterelor, în cazul hărților
procesate) fără să fiu nevoit să modific structura tabelelor. Tabelele se
creează automat la prima pornire, dacă nu există deja.

\section{Pipeline-ul Tier-2}
\label{sec:tier2}

Harta brută produsă de SLAM este, de obicei, destul de zgomotoasă: pereții
sunt cam zdrențuiți, apar puncte răzlețe și mici găuri. Pentru a obține ceva
care să semene cu un plan de etaj, am adăugat o etapă opțională de
post-procesare, pe care am numit-o ,,Tier-2" (Figura~\ref{fig:tier2}). Ea trece
harta printr-o succesiune de prelucrări clasice de imagine, descrise pe larg în
capitolul de teorie: curăță zgomotul prin operații morfologice, grupează
celulele în obiecte și aruncă grupurile prea mici (probabil zgomot), apoi
detectează pereții ca segmente drepte cu ajutorul transformatei Hough.
Rezultatul este o hartă mult mai curată, care arată ca o schiță de plan.

Am scris această prelucrare \textbf{numai cu NumPy}~\cite{harris2020numpy}, fără
biblioteci grele de viziune (OpenCV, scipy etc.). Motivul a fost să țin
instalarea de pe Pi cât mai ușoară; oricum, pe o hartă de dimensiuni obișnuite,
prelucrarea durează sub o secundă. Important este că harta procesată se salvează
separat, fără a o șterge pe cea originală — astfel, în interfața web pot comuta
oricând între varianta brută și cea curățată.

% Schema 7 — pipeline Tier 2
\begin{figure}[H]
\centering
\begin{tikzpicture}[
    node distance=0.3cm,
    every node/.style={font=\scriptsize},
    stage/.style={rectangle, draw, rounded corners, minimum width=5cm, minimum height=0.6cm, align=center},
    arr/.style={-{Latex[length=2mm]}}
]
    \node[stage, fill=blue!5]   (s1) {1. Filtru median 3$\times$3 (opțional)};
    \node[stage, fill=blue!10, below=of s1]  (s2) {2. Deschidere morfologică};
    \node[stage, fill=blue!15, below=of s2]  (s3) {3. Închidere morfologică};
    \node[stage, fill=blue!20, below=of s3]  (s4) {4. Etichetare componente conectate (BFS)};
    \node[stage, fill=blue!25, below=of s4]  (s5) {5. Transformata Hough probabilistică};

    \draw[arr] (s1) -- (s2);
    \draw[arr] (s2) -- (s3);
    \draw[arr] (s3) -- (s4);
    \draw[arr] (s4) -- (s5);

    \node[right=0.5cm of s5, font=\scriptsize] {$\to$ ProcessedMap};
\end{tikzpicture}
\caption{Pipeline-ul de post-procesare Tier-2 a hărților salvate (implementat în \texttt{recon\_db/\allowbreak postprocess.py}).}
\label{fig:tier2}
\end{figure}

\clearpage

% =============================================================================
% CAPITOLUL 10 — INTEGRARE & DEPLOYMENT
% =============================================================================
\chapter{Integrarea sistemului și desfășurarea}
\label{cap:deploy}

\section{Scriptul setup.sh}
\label{sec:setup-sh}

Sistemul este alcătuit din multe programe care trebuie pornite în ordinea
potrivită; ca să nu fac asta manual de fiecare dată, am scris un singur script
de pornire (\texttt{setup.sh}). El este și unicul mod recomandat de a porni
sistemul: o lansare manuală a nodurilor ar rata câțiva pași de pregătire a
mediului și ar duce la erori. Scriptul are mai multe \textbf{moduri}, în
funcție de ce vreau să fac: de la o pornire ușoară, doar cu interfața web și
date simulate (utilă pentru dezvoltat de pe laptop), până la modul complet, cu
toți senzorii reali, folosit pentru scanarea propriu-zisă; există și moduri
intermediare, pentru a testa separat IMU-ul sau LIDAR-ul.

Indiferent de mod, scriptul oprește mai întâi orice a rămas pornit de
data trecută, apoi verifică din timp că toate cele necesare sunt la locul lor
(Python, ROS2, Docker, portul senzorilor etc.) și se oprește cu un mesaj clar
dacă lipsește ceva. Fiecare program este pornit în propria fereastră a unei
sesiuni \textbf{tmux}, ca să-i pot urmări ușor mesajele. \textbf{Ordinea de
pornire contează}, pentru că nodurile depind unele de altele: puntea ESP32
trebuie să fie activă înainte ca celelalte să-i ceară date, EKF-ul înainte de
SLAM și așa mai departe, iar interfața web pornește ultima.

\section{Scriptul environment.sh}
\label{sec:environment-sh}

Dacă \texttt{setup.sh} \emph{pornește} sistemul, \texttt{environment.sh} îl
\emph{instalează}: pe un Raspberry~Pi proaspăt, el aduce și configurează tot ce
e nevoie — pachetele de sistem, ROS2, mediul Python, Docker, rețeaua WiFi. L-am
scris astfel încât să poată fi rulat de oricâte ori fără să strice ceva: înainte
de orice modificare verifică dacă nu cumva e deja făcută. Are și un mod de
\emph{verificare}, care doar bifează că totul e la locul lui pe un dispozitiv
deja configurat, fără să schimbe nimic.

Un detaliu care mi-a dat bătăi de cap a fost să pot porni interfața web pe
portul 80 (cel obișnuit pentru web) fără a rula totul ca administrator. Soluția
a fost să dau programului Python permisiunea specială necesară, dar aceasta
aduce, la rândul ei, o serie de mici complicații pe care scriptul le rezolvă
automat.

\section{Autostart prin systemd}
\label{sec:systemd}

Întrucât aparatul trebuie folosit „pe teren", fără tastatură și fără monitor,
trebuie ca, la simpla alimentare, totul să pornească singur. Pentru asta am
înregistrat sistemul ca un serviciu care rulează automat la pornirea
Raspberry~Pi~\cite{systemd}. Serviciul așteaptă întâi ca lucrurile de care
depinde să fie gata (Docker, rețeaua WiFi proprie) și chiar lasă câteva
secunde în plus pentru ca ESP32-ul să fie recunoscut de sistem după pornirea
la rece, abia apoi lansează stiva completă.

Am avut grijă și de cazurile în care ceva nu merge din prima: dacă pornirea
eșuează, serviciul reîncearcă de câteva ori, la interval de câteva secunde,
ceea ce rezolvă de la sine micile întârzieri de la pornirea la rece. Pornirea
automată poate fi, desigur, dezactivată, iar jurnalele pot fi urmărite în timp
real, lucruri utile atunci când vreau să depanez aparatul.

\clearpage

% =============================================================================
% CAPITOLUL 11 — TESTARE
% =============================================================================
\chapter{Testarea și validarea}
\label{cap:test}

\section{Strategia de testare}
\label{sec:test-strat}

Strategia de testare a proiectului este codificată ca regulă obligatorie
în \texttt{AGENT\_RULES.md}~§7 și se sprijină pe două principii. Întâi,
\textbf{fiecare nod are cel puțin un test} --- minim un test de tip
,,se încarcă fără să se prăbușească'', dar cu aserțiuni reale (un test care
trece întotdeauna nu este considerat test). În al doilea rând, testele
verifică \textbf{contractul, nu implementarea}: se validează că un nod
publică mesajele așteptate, pe topicele așteptate, cu tipurile așteptate,
fără a inspecta membri privați --- astfel încât o refactorizare internă să
nu spargă suita.

Suita este organizată pe două platforme: \textbf{pytest} pentru codul
Python (parser de cadre, integrator IMU, pipeline de post-procesare, modele
ORM, canale de date) și \textbf{gtest} prin \texttt{colcon~test} pentru
nodurile C++ (\texttt{draw\_node}, \texttt{sim\_sensor\_node}). Ambele suite
trebuie să fie ,,verzi'' înainte ca o etapă să fie declarată completă. La
ultima trecere de audit documentată, $59$ de cazuri pytest și $8$ de cazuri
gtest treceau, alături de verificarea construirii firmware-ului, a sintaxei
scripturilor shell și a construirii workspace-ului colcon.

\section{Teste unitare Python}
\label{sec:test-python}

\begin{table}[H]
\centering
\caption{Suita de teste Python (pytest).}
\label{tab:test-python}
\begin{tabular}{lcl}
\toprule
\textbf{Fișier} & \textbf{Cazuri} & \textbf{Acoperă} \\
\midrule
\texttt{test\_esp32\_uart\_bridge.py}   & 12 & parser framing (CRC, resync, erori) \\
\texttt{test\_imu\_yaw\_integrator.py}  & 12 & integrare unghi, cuaternion \\
\texttt{test\_postprocess.py}           & 20 & pipeline Tier-2 + Hough \\
\texttt{test\_db\_node.py}              & 9  & modele ORM, CRUD \\
\texttt{test\_recon\_webui.py}          & 6  & DataChannel, mock fallback \\
\bottomrule
\end{tabular}
\end{table}

\section{Teste unitare C++}
\label{sec:test-cpp}

Testele C++ sunt executate cu \texttt{colcon~test} și folosesc cadrul
gtest. \texttt{test\_draw\_node.cpp} (4 cazuri) verifică nodul
\texttt{draw\_node}: aranjarea grilei, operația de desenare (\textit{paint}),
ștergerea și limitarea (\textit{clamp}) a periei la marginile grilei.
\texttt{test\_sim\_sensor\_node.cpp} (4 cazuri) verifică generatorul de
LIDAR simulat: lovirea corectă a unui perete prin \textit{raycasting},
returnarea razei maxime în absența unui obstacol, conectivitatea camerei
generate (verificată prin BFS) și blocarea razei de către un obstacol
interior. Aceste teste permit exersarea pipeline-ului SLAM~$\to$~bază de
date~$\to$~interfață web fără hardware fizic, prin nodul de simulare.

\section{Validarea pe hardware real}
\label{sec:hardware-val}

Pe lângă testele unitare, sistemul a fost validat pe hardware real, cu
rezultatele documentate în \texttt{STATUS.md}~§6 și
\texttt{REMAINING\_ISSUES.md}. La pornirea în modul \texttt{full}, apăsarea
butonului \textit{Start scan} pornește motorul LIDAR, după care
\texttt{/scan} este publicat la $\sim\SI{6}{\hertz}$ și harta
\texttt{/map} este actualizată la $\SI{1}{\hertz}$ în mod susținut;
apăsarea \textit{Pause} oprește motorul și îngheață harta curat.
Inventarul de topice ROS confirmă ratele așteptate
($\texttt{/tf}\approx\SI{74}{\hertz}$, $\texttt{/odom}\approx\SI{24}{\hertz}$,
$\texttt{/imu/data}=\SI{100}{\hertz}$,
$\texttt{/esp32/diagnostics}=\SI{1}{\hertz}$), iar funcția
\textit{Clear map} resetează SLAM-ul în $\sim\SI{40}{\milli\second}$ (în
pauză) / $\sim\SI{300}{\milli\second}$ (activ).

Validarea legăturii UART a fost realizată prin captura a peste $3500$ de
cadre IMU și $33$ de \textit{heartbeat}-uri pe o fereastră de $33$~secunde,
\textbf{fără nicio eroare CRC}. Întregul lanț a fost, de asemenea,
verificat end-to-end pe banc (rotirea LIDAR-ului, generarea scanărilor prin
calea PTY, popularea hărții). Testul de plimbare propriu-zis într-o cameră
de $5\times\SI{5}{\meter}$, cu măsurarea acurateței dimensionale a pereților
față de criteriile de acceptanță din §\ref{sec:obiective}, este planificat
ca etapă viitoare~H4.

\clearpage

% =============================================================================
% CAPITOLUL 12 — CONCLUZII
% =============================================================================
\chapter{Concluzii}
\label{cap:concluzii}

\section{Sumarul realizărilor}
\label{sec:summary}

La final, sistemul funcționează cap-coadă. Senzorul LD14P, citit prin ESP32 și
redat pe Pi prin pseudo-terminal, alimentează slam\_toolbox, care construiește
în timp real harta de ocupare; aceasta poate fi urmărită în interfața web și
salvată în PostgreSQL. Orientarea aparatului vine din integrarea giroscopică
trecută prin EKF, iar SLAM-ul îi corectează permanent deriva. În plus, hărțile
salvate pot fi curățate ulterior, ca să arate ca un plan de etaj, nu ca o
imagine zgomotoasă.

Cerințele practice pe care mi le propusesem la început au fost îndeplinite.
Costul total al componentelor (în jur de $\SI{166}{\euro}$) a rămas sub pragul
de $\SI{200}{\euro}$; dispozitivul este complet autonom, datorită propriei
rețele WiFi; poate fi pornit de oricine, dintr-un singur buton; iar pornirea
automată la alimentare îl face utilizabil fără tastatură și fără monitor.
Testele pe hardware au arătat un flux de date stabil și o hartă care se
actualizează constant. Singurul criteriu pe care nu l-am putut încă măsura
riguros este acuratețea dimensională (eroarea pe pereții lungi): acesta rămâne
de verificat printr-un test de plimbare într-o cameră reală, planificat ca pas
următor.

\section{Contribuții originale}
\label{sec:contributions}

Dincolo de integrarea unor componente open-source consacrate, lucrarea
aduce următoarele contribuții originale:

\begin{enumerate}\itemsep0pt
    \item \textbf{Calea LIDAR prin ESP32 cu pseudo-terminal (PTY)} ---
          soluție neconvențională pentru transmiterea fluxului de date
          LIDAR către driverul ROS2 fără a folosi UART-ul nativ al Pi-ului,
          unificând toate datele senzoriale pe o singură legătură USB-CDC;
    \item \textbf{Patch-ul \texttt{42688f6}} pe driverul
          \texttt{ldlidar\_stl\_ros2} pentru fixarea geometriei la 720 de
          fascicule, descoperit prin investigarea în profunzime a bug-ului
          de blocare silențioasă a actualizării hărții cauzat de blocarea
          numărului de fascicule în Karto;
    \item \textbf{Documentarea explicită} a problemei culorilor de fir
          non-standard ale LD14P, un \textit{gotcha} cu valoare practică
          ridicată pentru orice replicare a montajului;
    \item \textbf{Pipeline de post-procesare pur-NumPy} (filtru median,
          morfologie, componente conectate, transformata Hough
          probabilistică) fără dependențe grele de tip scipy/OpenCV,
          potrivit pentru o instalare lejeră pe Pi;
    \item \textbf{Integrarea WiFi AP + DHCP + DNS local} astfel încât
          dispozitivul să fie complet auto-conținut în teren, fără a depinde
          de o rețea preexistentă.
\end{enumerate}

\section{Limitări cunoscute}
\label{sec:limitations}

Sistemul are limitări cunoscute, documentate transparent în
\texttt{REMAINING\_ISSUES.md} și \texttt{STATUS.md}~§7:

\begin{itemize}\itemsep0pt
    \item \textbf{Lipsa autentificării} pe interfața web --- acceptabilă în
          modelul de operare pe AP izolat, dar care necesită autentificare
          de bază înainte de orice expunere pe un LAN partajat;
    \item \textbf{EKF doar pe orientare} --- translația nu provine din
          dead-reckoning inerțial (din cauza biasului \texttt{ZA\_OFFSET}),
          ci exclusiv din corecția \texttt{map}~$\to$~\texttt{odom} a SLAM-ului;
    \item \textbf{Drift giroscopic} la dispozitivul staționar
          ($\sim\SI{0.7}{\degree\per\second}$), tolerabil deoarece este
          corectat de potrivirea de scanări, dar prezent;
    \item \textbf{Acoperire de test incompletă} pentru opcode-urile
          \texttt{LIDAR\_FRAME}/\texttt{LIDAR\_EN}/\texttt{LIDAR\_ACK},
          validate deocamdată doar live, nu și prin teste unitare;
    \item \textbf{Driverul vendorat ca \textit{gitlink} fără submodul} ---
          o clonare naivă pierde patch-urile locale (vezi
          §\ref{sec:ldlidar-driver}), necesitând re-vendorarea explicită;
    \item \textbf{Fragilitatea modelului eventlet~+~rclpy} --- invariantele
          de \textit{threading} sunt obligatorii și ușor de încălcat de cod
          nou (§\ref{sec:threading}).
\end{itemize}

\section{Direcții viitoare}
\label{sec:future}

Direcțiile viitoare urmează etapele rămase din foaia de parcurs a
proiectului:

\begin{itemize}\itemsep0pt
    \item \textbf{H3.1}: calibrarea de banc a IMU-ului și posibila adoptare
          a filtrului Madgwick~\cite{madgwick2014} pentru
          \textit{roll}/\textit{pitch};
    \item \textbf{H4}: primul test real de plimbare într-o cameră
          necunoscută, cu cuantificarea acurateței dimensionale;
    \item \textbf{H5}: o mașină de stare reală a sesiunii de scanare
          (IDLE~$\to$~SCAN~$\to$~PAUSED~$\to$~SAVED), care să controleze
          scrierile în baza de date;
    \item \textbf{H6}: încastrarea (\textit{enclosure}) imprimată 3D și
          integrarea fizică finală (fixarea acumulatorului portabil de
          22,5 W în carcasă), plus oprirea soft prin butonul SHUTDOWN;
    \item \textbf{extensii speculative}: cartografiere 3D cu un servo de
          înclinare a LIDAR-ului și export în format STL/OBJ pentru
          aplicații AR/VR.
\end{itemize}

\clearpage

% =============================================================================
% BIBLIOGRAFIE
% -----------------------------------------------------------------------------
% Capitolul + intrarea în Cuprins sunt generate automat de \bibliography{...}
% combinat cu pachetul tocbibind și redefinirea lui \bibname din preambul.
% Stilul IEEEtran sortează implicit referințele în ordinea citării în text
% (\cite{...}), conform cerințelor UPT.
% =============================================================================
\bibliography{references}

% =============================================================================
% DECLARAȚIE DE AUTENTICITATE (ultimă pagină)
% =============================================================================
\clearpage
\thispagestyle{coperta}

\begin{center}
{\fontsize{14pt}{17pt}\bfseries\MakeUppercase{Declarație de autenticitate}}\\[0.4cm]
{\fontsize{12pt}{14pt}\itshape (a lucrării de finalizare a studiilor)}
\end{center}

\vspace{0.8cm}

\noindent Subsemnatul/a, \candidatname{}, candidat la examenul de
finalizare a studiilor la Facultatea de Automatică și Calculatoare,
specializarea \specializare{}, sesiunea \sessionname{}, declar pe propria
răspundere că lucrarea de finalizare a studiilor cu titlul:

\medskip

\noindent\textbf{,,\workTitle{}''}

\medskip

\noindent este elaborată de mine și nu a mai fost prezentată niciodată la
o altă facultate sau instituție de învățământ superior din țară sau din
străinătate.

\medskip

\noindent De asemenea, declar că toate sursele utilizate, inclusiv cele de
pe Internet, sunt indicate în lucrare, cu respectarea regulilor de evitare
a plagiatului:
\begin{itemize}\itemsep0pt
    \item toate fragmentele de text reproduse exact, chiar și în
          traducere proprie din altă limbă, sunt scrise între ghilimele
          și dețin referința precisă a sursei;
    \item reformularea în cuvinte proprii a textelor scrise de către alți
          autori deține referința precisă;
    \item codul sursă, imaginile etc.\ preluate din proiecte open-source
          sau alte surse sunt utilizate cu respectarea drepturilor de
          autor și dețin referințe precise;
    \item rezumarea ideilor altor autori deține referința precisă la
          textul original.
\end{itemize}

\vspace{1.5cm}

\noindent
\begin{minipage}[t]{0.5\textwidth}
    Timișoara,\\
    Data: \rule{4cm}{0.4pt}
\end{minipage}%
\begin{minipage}[t]{0.5\textwidth}
    Semnătura candidatului\\[1.0cm]
    \rule{5cm}{0.4pt}
\end{minipage}

\end{document}
